\documentclass[fontset = fandol, svgnames]{ctexrep}

\ctexset{
  chapter        = {
    format       = \huge\bfseries\raggedright,
    number       = \arabic{chapter},
    name         = {},
    tocline      = \CTEXifname{\protect\numberline{\thechapter}}{}#2,
  },
  section/format = \Large\bfseries\raggedright,
}

\usepackage[a4paper, margin = 2.4cm]{geometry}
\usepackage{booktabs, array}
\usepackage{fontawesome7, fancyqr, hologo}
\usepackage[colorlinks,
  linkcolor  = blue,
  citecolor  = blue,
  urlcolor   = blue,
  pdfpagelayout = SinglePage,
  bookmarksnumbered]{hyperref}
\hypersetup{
  pdftitle   = {和AI好好说话：零基础提示工程指南},
  pdfauthor  = {吕峰军（明城）},
  pdfsubject = {提示工程入门教程},
  pdfkeywords = {提示工程, Prompt Engineering, 提示词, 大语言模型, LLM, AI, ChatGPT, Zero-shot, Few-shot, Chain of Thought, 多模态, RAG},
  pdfcreator = {XeLaTeX},
  pdflang    = {zh-CN},
}

\usepackage{tikz}
\usetikzlibrary{shapes.geometric, calc}
\tikzset{
  score stars/.style = {
    % shape
    star, star points = 5, star point ratio = 2.25, scale = .8,
    % color 
    draw = gray, fill = #1,
    % others
    inner sep = 0.14em, anchor = outer point 3
  }
}
\newcommand\stars[1]{%
  \begin{tikzpicture}
    % Draw five stars. For #1 = "2.3", fill the 1st and 2nd stars as gray,
    % and fill the 3rd to 5th stars as white.
    \foreach \i in {1, ..., 5} {
      \pgfmathsetmacro\starcolor{\i<=#1 ? "gray" : "white"}
      \node [score stars = \starcolor] (star\i) at (\i*0.8em, 0) {};
    }
    % For #1 = "2.3", let \partstar = "3" and \starpart = "0.3".
    % Then fill the left 30% part of the 3rd star as gray after clipping.
    \pgfmathsetmacro\partstar{#1>int(#1) ? int(#1+1) : 0}
    \ifnum\partstar>0
      \pgfmathsetmacro\starpart{#1-(int(#1))}
      \coordinate (upper left)
        at (star\partstar.outer point 2 |- star\partstar.outer point 1);
      \coordinate (upper right)
        at (star\partstar.outer point 5 |- star\partstar.outer point 1);
      \coordinate (lower right)
        at (star\partstar.outer point 5 |- star\partstar.outer point 4);
      \clip (upper left) rectangle
            ({$ (upper left)!\starpart!(upper right) $} |- lower right);
      \node [score stars = gray] at (\partstar*0.8em, 0) {};
    \fi
  \end{tikzpicture}%
}
\newdimen\ttwd \newbox \xttwdbox
\setbox\xttwdbox\hbox{\normalsize\ttfamily x}
\ttwd=\wd\xttwdbox

\usepackage[os = win]{menukeys}
\renewmenumacro{\menu}[>]{angularmenus}
\renewmenumacro{\keys}[+]{shadowedroundedkeys}
\renewcommand\RSsmallest{5pt}
\protect\renewcommand\faWindows{%
  \tikz[rounded corners = .1pt, baseline = -.25em]
    \foreach \i in {0, 90, 180, 270}
      \fill [ rotate = \i ] (-.4em, .4em) rectangle (-.025em, .025em);
}

\usepackage{listings, fancyvrb}
\usepackage{tcolorbox}
\tcbuselibrary{breakable, skins}

% 提示词代码块样式（用户输入的提示词）
\newtcolorbox{promptbox}{%
  enhanced,
  breakable,
  colback    = blue!3!white,
  colframe   = blue!45!black,
  leftrule   = 0.4pt,
  rightrule  = 0.4pt,
  toprule    = 0.4pt,
  bottomrule = 0.4pt,
  arc        = 0pt,
  left       = 10pt,
  right      = 10pt,
  top        = 8pt,
  bottom     = 8pt,
  fontupper  = \small,
  before upper = {\setlength{\parskip}{0.6em}\setlength{\parindent}{0pt}},
}

% AI 回答样式
\newtcolorbox{aibox}{%
  enhanced,
  breakable,
  colback    = green!3!white,
  colframe   = green!40!black,
  leftrule   = 0.4pt,
  rightrule  = 0.4pt,
  toprule    = 0.4pt,
  bottomrule = 0.4pt,
  arc        = 0pt,
  left       = 10pt,
  right      = 10pt,
  top        = 8pt,
  bottom     = 8pt,
  fontupper  = \small,
  before upper = {\setlength{\parskip}{0.6em}\setlength{\parindent}{0pt}},
}

% 小贴士样式
\newtcolorbox{tipbox}{%
  enhanced,
  breakable,
  colback    = orange!3!white,
  colframe   = orange!50!black,
  leftrule   = 0.4pt,
  rightrule  = 0.4pt,
  toprule    = 0.4pt,
  bottomrule = 0.4pt,
  arc        = 0pt,
  left       = 10pt,
  right      = 10pt,
  top        = 8pt,
  bottom     = 8pt,
  fontupper  = \small,
  before upper = {\setlength{\parskip}{0.6em}\setlength{\parindent}{0pt}},
}

\lstset{
  breaklines      = true,
  columns         = fullflexible,
  showstringspaces= false,
  tabsize         = 4,
  gobble          = 1,
  numbers         = left,
  numberstyle     = \tiny\ttfamily,
  numbersep       = \ccwd,
  frame           = lines,
  rulecolor       = \color{blue!40},
  backgroundcolor = \color{lightgray!20},
  language        = bash,
  alsoletter      = -,
  basicstyle      = \small\ttfamily,
  commentstyle    = \color{olive!80!black},
  stringstyle     = \color{brown},
  emph            = { bash, cd, cp, dpkg, find, grep, ln, md5sum, mkdir,
                      mount, rm, rmdir, set, sha512sum, sudo, umount },
  emphstyle       = \color{blue}\bfseries,
  morekeywords    = { add-apt-repository, apt, brew, certutil, evince,
                      fc-cache, fc-list, fmtutil, fmtutil-user, fmtutil-sys,
                      gedit, kpsewhich, l3build, latexmk, mkfontdir,
                      mkfontscale, notepad, tex, texdoc, tlmgr, pdflatex,
                      visudo, wget, wsl, xelatex },
  keywordstyle    = \color{teal}\bfseries
}
\lstdefinelanguage{json}{
  basicstyle   = \small\ttfamily\color{DarkBlue},
  numberstyle  = \tiny\ttfamily\color{black},
  alsoletter   = ",
  morecomment  = [l]{//},
  morekeywords = {"latexmk", "latexmkpdf", "latexmkxe", "never", "tab", "external"},
  keywordstyle = \color{FireBrick}\bfseries,
  literate     ={*[{\textcolor{DarkGreen}{[ }}1  ]{\textcolor{DarkGreen}{ ]}}1
                 \{{\textcolor{DarkRed}  {\{}}1 \}{\textcolor{DarkRed}  {\}}}1
                  :{\textcolor{black}    {:}} 1  ,{\textcolor{black}    {,}} 1}
}
\lstdefinelanguage{mwe}{
  language     = {[LaTeX]TeX},
  texcsstyle   = *\color{violet},
  morekeywords = {article, ctexart, document},
  literate     =*\{{\textcolor{magenta}\{}1 \}{\textcolor{magenta}\}}1
}
\lstMakeShortInline [ breaklines = true, basicstyle = \ttfamily,
                      emphstyle = {}, keywordstyle = {} ] {"}
\VerbatimFootnotes

\title{\bfseries 和AI好好说话：零基础提示工程指南%
  \thanks{\url{https://github.com/mingcheng/prompt-engineering-guide-for-beginners}}%
}
\author{吕峰军（明城）%
  \thanks{\href{mailto:mingcheng@outlook.com}%
    {\texttt{mingcheng@outlook.com}}}%
}
\date{\today}

\begin{document}

\maketitle

%!%
% Copyright (c) 2026 mingcheng <mingcheng@outlook.com>
% 
% File: preface.tex
% Author: mingcheng <mingcheng@outlook.com>
% File Created: 2026-03-29 00:58:27
% 
% Modified By: mingcheng <mingcheng@outlook.com>
% Last Modified: 2026-03-29 07:14:47
%%

% !TeX root = ../prompt-engineering-guide-for-beginners.tex

\chapter*{前言}

\section*{关于作者}

大家好，我是吕峰军（明城），一个写了十几年代码的技术人\footnote{\href{https://github.com/mingcheng}{我的 GitHub 主页}}。
平时在互联网公司做后端和架构相关的工作，同时也是一个多年的开源爱好者。
这几年 AI 的发展太快了，我也一直在关注和尝试各种 AI 工具。从最初的新奇到日常的依赖，过程中踩了不少坑，也攒下了一些经验。

\section*{为什么写这本小册子？}

写这本小册子的起因特别朴素。

公司里越来越多的小伙伴开始用 AI 了，但大家的使用效果参差不齐：有人用得飞起，三分钟搞定一篇周报；有人折腾半天，AI 给出来的东西完全不能用。

问题出在哪儿？大多数时候，不是 AI 不行，而是我们跟它「说话的方式」不对。
同样一个需求，换一种问法，结果可能天差地别。
这就是提示工程（Prompt Engineering）的价值所在：学会跟 AI 好好说话，让它真正成为你的得力助手。

这本小册子最初源自团队内部的分享。后来有同事说，身边很多朋友都有同样的困惑，建议我整理得更详细些，分享给更多人。
于是我趁着业余时间不断补充和完善，最终形成了你现在看到的这本小册子。

不管你是刚接触 AI 的纯新手，还是已经用了一阵但总觉得「差点意思」的老用户，我都希望这本小册子能帮到你。

\section*{怎么读这本小册子？}

「提示工程」听起来有点高大上，但它的底层逻辑并不玄乎。本质上它是用自然语言作为代码，对概率庞大的神经网络模型进行一次“黑盒编程”。
这不是什么需要门槛的前沿技术，而是一项极为实用的逻辑外脑驾驭技能。你不需要熟练掌握任何编程语言，不需要弄懂 Transformer 的底层结构，只要会打字、拥有清晰的结构化思维就能学会。

当然，掌握它也需要一些时间和逻辑训练。刚开始可能会觉得有点抽象甚至是“玄学”，但别担心，书里的内容摒弃了艰涩的算法推导，尽量用通俗易懂的工程化类比，并配合了大量的真实示例和可复用的模板，帮你快速上手这门“赛博手艺”。

在开始之前，我想给你三个建议：

\textbf{第一，别光看，要动手。}每一章里的例子和模板，都建议你打开 AI 工具实际试一遍。看十遍不如自己写一遍，亲手体验一次胜过读十页理论。

\textbf{第二，别追求完美，先用起来。}提示词没有标准答案，同一个需求可以有很多种写法。先从简单的开始，慢慢迭代优化，你会越用越顺手。

\textbf{第三，保持好奇心，但不必焦虑。}AI 模型和提示技巧都在飞速演进，你不需要追逐每一个新概念。只要掌握了核心原则，新工具出来时你也能很快上手。

\section*{关于更新与贡献}

AI 领域变化太快了，所以这本小册子也不会一成不变。我会持续更新内容，跟上最新的模型能力和技术发展。

这本小册子在 GitHub 上完全开源（基于 Apache 2.0 许可证）。
如果你发现了错误、有更好的案例、或者想补充新的内容，非常欢迎通过 Pull Request 来贡献。
同样，如果你有任何意见和建议，也欢迎提 Issue 告诉我。

众人拾柴火焰高，期待和你一起把这份资料做得更好。

\section*{致谢}

感谢公司同事们在日常工作中的实践和反馈，书中很多真实案例都来自大家的使用经验。
也感谢每一位愿意花时间阅读这本小册子的你。

希望它能帮你在 AI 时代多一项趁手的技能。那我们开始吧！

\section*{关键词}

提示工程（Prompt Engineering）、提示词（Prompt）、大语言模型（LLM）、人工智能（AI）、ChatGPT、零样本提示（Zero-shot）、少样本提示（Few-shot）、思维链（Chain of Thought）、多模态（Multimodal）、检索增强生成（RAG）


\tableofcontents

%!%
% Copyright (c) 2026 mingcheng <mingcheng@outlook.com>
% 
% File: ch01-icebreaker.tex
% Author: mingcheng <mingcheng@outlook.com>
% File Created: 2026-03-27 14:57:13
% 
% Modified By: mingcheng <mingcheng@outlook.com>
% Last Modified: 2026-03-29 07:14:37
%%

% !TeX root = ../prompt-engineering-guide-for-beginners.tex

\chapter{破冰：把 AI 当成你的超级实习生}
\label{ch:icebreaker}

如果你觉得 AI 离自己很远，甚至有些高深莫测，那这一章就是专门为你准备的。
读完之后你会发现，AI 既不神秘，也不可怕。
它本质上就是一个特别能写字的助手，而关键在于你怎么跟它「说话」。

\section{什么是提示工程（Prompt Engineering）？}

简单来说，提示工程就是「面向大型概率模型进行自然语言编程」的学问。

你在 AI 对话框里输入的每一段文字，都叫做「提示词」（Prompt）。而提示工程研究的，就是如何写出更好的提示词，这实际上是在为模型设定一个精密的“初始上下文”，从而引导并在数学概率上极大限制它生成你想要的结果。你可以把它理解为新时代的「黑盒编程语言」：只不过你用的不再是需要编译和严谨语法的 Python 或 C++，而是我们日常交流、具有一定歧义的自然语言。

打个比方：去餐厅点菜时，如果只说「来个菜」，服务员完全不知道你想吃什么；但如果你说「来一份不辣的番茄炒蛋，少油少盐」，厨师就能精准地做出你想要的菜。跟 AI 打交道，道理完全一样。

提示工程不要求你会写代码，也不要求你懂机器学习算法。它是一项\textbf{人人都能学会的沟通技能}，只要你会打字、能表达需求，就可以上手。

\begin{tipbox}
很多新手在初次使用 AI 时，会担心自己的表达不够书面化或有语病，实际上大模型的语义理解能力非常强，对语言的容错率极高。就像与真人助理交代工作一样用大白话提问即可，完全不必过分纠结语法与辞藻。
\end{tipbox}

\section{大语言模型到底是个啥？}

你肯定听过 ChatGPT、Claude、Gemini 这些名字。它们背后的核心技术叫做「大语言模型」（Large Language Model，简称 LLM）。

不用被术语吓到。从技术底层来看，你可以把大语言模型想象成一个基于 Transformer 架构的「超级文字接龙机器」：它在训练阶段阅读过数十 TB 的乃至更庞大的互联网数据（包括文章、书籍和代码库），通过自注意力机制（Self-Attention）死记硬背下了人类语言中词与词之间的高维统计学规律。当你给它提供一句话时，它其实是在计算概率分布，预测接下来最可能出现的下一个词（Token），如此循环往复，一个词一个词地「接」出一段完整的回答。比如，当你输入「白日依山」，它就会自然而然地算出「尽」的出现概率最高（可能高达 99\%），从而接上这个字。

不仅如此，现代的对话大模型（如 ChatGPT）还经过了“指令微调（SFT）”和“人类反馈强化学习（RLHF）”。如果是原始的“基础模型”，它只会单纯接龙；但在经过指令微调后，它才真正变成了一个能听懂人类指令并“服从命令”的 AI 助手。

换句话说，AI 并不是真的像人类一样在利用物理大脑「理解」你说了什么，也不是在用独立意识「思考」问题的答案。它的本质是在做一个极其复杂的数学概率运算，生成一段在人类看来连贯且符合常理的响应。

这意味着两件事：

\begin{itemize}
  \item 它能写出非常流畅、专业的内容，因为它见过的好文章足够多；
  \item 它也可能一本正经地胡说八道，因为它只在意「说得通」，并不在意「是不是真的」。
\end{itemize}

理解这一点很重要，因为它直接决定了我们应该怎么使用 AI、怎么看待它给出的回答。

\section{为什么你的 AI 看起来很笨？}

很多人第一次用 AI 时，问了一个问题，得到一个不太满意的回答，就下结论说「AI 不行」。

但真相往往是：不是 AI 不行，而是我们提问的方式有问题。

来看一个对比：

\begin{itemize}
  \item \textbf{模糊的提问}：「帮我写个东西。」
  \item \textbf{清晰的提问}：「帮我写一封给客户的邮件，告诉对方项目进度延迟了两周，语气要礼貌但诚恳，控制在 200 字以内。」
\end{itemize}

第一种问法，AI 根本不知道你想要什么，只能瞎猜，给出的内容往往大而无当。第二种问法，AI 有了明确的上下文（特定的对象、背景、语气和字数限制），自然能给出靠谱得多的回答。

这就是提示工程中的一条黄金法则：\textbf{垃圾输入 = 垃圾输出}（Garbage In, Garbage Out）。你给 AI 的指令越模糊，回答就越空洞；指令越具体，回答就越精准。为了让 AI 表现得像个专业的得力助手，你首先得给它提供清晰、充分的背景信息。

再举一个日常例子。假设你想让 AI 帮你策划一顿朋友生日聚餐：

\begin{itemize}
  \item \textbf{模糊版}：「帮我推荐一家餐厅。」
  \item \textbf{清晰版}：「帮我推荐一家上海静安区适合 6 人聚餐的餐厅，预算人均 200 元左右，朋友中有一个人吃素，最好有独立包间，氛围适合过生日。」
\end{itemize}

前者，AI 可能给你推荐从街边麻辣烫到高端法餐的一串清单，什么都有但什么都没用。后者，AI 就能在地段、人数、预算、饮食限制等多项约束条件的交叉筛选下给出靠谱得多的答案。

好消息是，学会「好好提问」并不难，接下来的章节会一步步带你掌握。

\begin{tipbox}
如果 AI 给出的一次回答不尽如人意，不要立刻开启新对话或认为求助失败。你应该直接在对话框里追问和修正，例如补充一句「总结得太长了，帮我缩短到 100 字」或「语气再活泼一些」。这种不断微调的「多轮对话」，本身就是提示工程中非常高频且关键的操作。
\end{tipbox}

\section{市面上的模型怎么选？}

目前市面上有很多 AI 工具可供选择，各有特点。以下是一些主流选项的简要介绍：

\begin{itemize}
  \item \href{https://chatgpt.com}{\textbf{ChatGPT}}（OpenAI）：目前最知名的 AI 对话工具，功能全面，支持文本、图片、代码等多种任务，有免费版和付费版。
  \item \href{https://claude.ai}{\textbf{Claude}}（Anthropic）：擅长长文本处理和细致的分析任务，语气相对温和严谨，适合需要高质量文字输出的场景。
  \item \href{https://gemini.google.com}{\textbf{Gemini}}（Google）：谷歌推出的 AI 助手，和 Google 生态深度整合，适合习惯使用谷歌全家桶的用户。
  \item \href{https://qwen.ai}{\textbf{Qwen}}（阿里巴巴）: 阿里巴巴的 AI 模型，支持多语言输入，适合需要处理中文和其他语言混合内容的用户。
  \item \href{https://kimi.moonshot.cn}{\textbf{Kimi}}（月之暗面）：国产 AI 工具，支持超长文本输入，适合需要阅读和分析大量中文资料的用户。
  \item \href{https://www.deepseek.com}{\textbf{DeepSeek}}：国产开源模型的代表，性价比高，在推理和代码方面表现不俗。
  \item \href{https://www.zhipuai.cn}{\textbf{智谱 GLM}}：清华背景的国产大模型，提供丰富的 API 接口，适合有开发需求的用户。
\end{itemize}

如果你是纯新手，建议从 ChatGPT、Qwen、Kimi 或 DeepSeek 中任选一个开始。它们注册简单、界面友好，免费额度足够入门练习。不必纠结哪个「最好」，对于日常使用来说，\textbf{哪个顺手就用哪个}。随着你对 AI 了解的不断加深，自然会体会到不同模型在特定任务上的细微优势。

值得一提的是，目前很多平台都提供了「深度思考」或「推理模型」的选项（如 OpenAI 的 o1/o3、DeepSeek R1 等）。这类模型在回答前会先进行一段内部推理，适合处理数学、逻辑和复杂分析等任务。入门阶段用普通的对话模型就足够了，关于推理模型的更多细节将在第~\ref{sec:reasoning-models}~节介绍。

顺便推荐一下 \href{https://openrouter.ai}{OpenRouter}、\href{https://zenmux.ai}{ZenMux} 等聚合平台，它们提供统一的接口服务，让你可以在不同模型之间自由切换和对比试用，方便找到最适合自己的那个。
此外，\href{https://cherry-ai.com}{Cherry Studio}、\href{https://anythingllm.com}{AnythingLLM} 等第三方客户端工具也值得关注，它们不仅提供了多模型切换的便利，还内置了提示词管理功能，能帮助你更高效地组织与沉淀自己的提示词库，是进阶使用者的好帮手。

我个人使用 Claude 和 Gemini 比较多，但说实话 Claude 的费用偏高。如果你只是想练习提示词写作，完全可以先用 Qwen 或 Kimi 等有免费额度的模型，等熟练之后再考虑升级到更强大的付费模型。

\begin{tipbox}
无论选择哪个模型，在工作中使用 AI 时都请务必注意数据安全与隐私。请勿将公司的核心商业机密、内部代码库的密钥或客户的敏感个人信息毫无处理地直接输入给公共 AI 服务。
\end{tipbox}

\section{AI 的能力边界：它能做什么、不能做什么}

在深入学习之前，有必要先了解 AI 的能力边界，这样可以少走弯路。

\subsection*{AI 擅长的事情}

AI 最大的优势在于处理文字相关的任务：速度快、不怕重复、不会抱怨加班。具体来说，它擅长以下几类工作：

\begin{itemize}
  \item \textbf{文本生成}：写邮件、写文案、写报告、写故事。给一个方向，它就能快速产出初稿，省去「从零开始」的痛苦。
  \item \textbf{翻译}：中英互译乃至多语言翻译。日常沟通级别的翻译，AI 已经足够好用。
  \item \textbf{总结和提炼}：帮你把长文章浓缩成几句话。面对一份 50 页的报告，让 AI 先帮你抓重点，效率提升显著。
  \item \textbf{头脑风暴}：提供创意和灵感。当你卡在一个想法上不知如何展开时，AI 是个不错的「点子发生器」。
  \item \textbf{格式转换}：把文字整理成表格、列表、JSON 等格式。这类机械性的排版工作，交给 AI 再合适不过。
  \item \textbf{代码辅助}：写代码、改 Bug、解释代码逻辑。即使你不是程序员，AI 也能帮你写一些简单的自动化脚本。
\end{itemize}

\begin{tipbox}
遇到复杂任务时，不要指望 AI 能「一次写出完美成品」。最好的心态是把它当成一个不知疲倦的「草稿生成器」。由它快速完成从 0 到 80 分的草图与体力活，你再负责从 80 到 100 分的精修与灵魂注入，这才是目前最高效的人机协作模式。
\end{tipbox}

\subsection*{AI 不擅长的事情}

与此同时，AI 也有明显的短板。认清这些不足，能帮你避免在错误的场景中使用它：

\begin{itemize}
  \item \textbf{实时信息}：如果没有联网功能，AI 不知道今天的新闻或最新的股价。它的知识有一个「截止日期」，此后的事情一概不知。
  \item \textbf{精确计算}：复杂的数学运算它可能算错。AI 不是计算器，它是在「预测下一个最可能的词」，而不是在「算数」。虽然新版模型在这方面持续改进（例如能够自动调用代码来辅助计算），但涉及重要财务或业务数据时，仍建议自行验算。
  \item \textbf{事实准确性}：它可能会「编造」听起来很真实的假信息（这在一定程度上被称为「幻觉」，Hallucination）。比如你问它某本冷门书籍的作者，它可能自信满满地报出一个错误的名字。因此，在查阅专业文献或获取硬核知识时，务必进行交叉验证。
  \item \textbf{读心术}：你没说的东西，它猜不到。AI 只能根据你提供的信息来工作，无法感知你的真实意图、工作背景或情绪状态。
\end{itemize}

值得注意的是，虽然越来越多的模型已经支持联网搜索、调用外部工具等扩展功能（例如 MCP，即 Model Context Protocol，一种让 AI 连接外部数据源的标准化协议），但这并不意味着它就能完全克服上述短板。你仍然需要对 AI 的回答保持审慎态度，尤其是涉及重要决策或敏感信息时。

\subsection*{技术视角的常见误解}

最后，澄清几个对于 LLM 底层运作常见的「认知误区」：

\begin{itemize}
  \item \textbf{大模型不是搜索引擎}。搜索引擎是去互联网倒排索引中查找已有的映射网页，而 AI 则是根据学到的语言分布模式凭空“计算”并生成新的文字。它不是在特定的硬盘上「查」资料，而是在神经网络权重中「推演」回答。
  \item \textbf{大模型不是关系型数据库}。它没有内部的一个 SQL 乃至图数据库供你完美精确的提取信息。它的知识是模糊地压缩在千亿个神经元参数里的。因此绝对不能把它当成严格的词典或百科全书来信任。
  \item \textbf{无状态（Stateless）本质：AI 默认不具有“长久记忆”}。作为纯粹的推理引擎，每次你开启新对话，它的计算状态都会被清零。它就像一个「永远清醒但频繁失忆」的打字机，对你上个窗口说过的话一无所知（除非平台本身开发了额外的记忆检索外挂体系并自动组装到了你的系统提示词里）。
\end{itemize}

从合理的期待出发，逐步探索 AI 的能力边界，你会发现它是一个非常强大的生产力工具，能在很多场景下帮你节省大量时间和精力。
了解了这些能力边界，你就能把 AI 放在最擅长的位置上使用，既不会因期望过高而失望，也不会因误解而低估它的价值。

理论已经铺设完毕，接下来最好的学习方式就是亲自上手。下一章，我们会在五分钟内完成你的第一次 AI 对话。

% !TeX root = ../prompt-engineering-guide-for-beginners.tex

%%%%%%%%%%%%%%%%%%%%%%%%%%%%%%%%%%
%% 第二章：动手试试
%%%%%%%%%%%%%%%%%%%%%%%%%%%%%%%%%%
\chapter{动手试试：五分钟开启你的第一次 AI 对话}
\label{ch:first-try}

上一章我们聊了 AI 是什么、不是什么。这一章，咱们直接动手。
读再多理论，不如自己试一次来得印象深刻。
五分钟之后，你就会完成人生中第一次（或者更有质量的一次）AI 对话。

\section{选一个工具，注册并打开}

首先，你需要选一个 AI 对话工具。如果你尚未注册过任何 AI 工具，推荐从以下两个对新手最友好的国产模型开始练习。它们的注册流程比海外产品简单，且交互界面十分友好：

\begin{itemize}
  \item \href{https://qwen.ai}{\textbf{通义千问 (Qwen)}}：由阿里巴巴推出，支持直接使用手机号注册，免费额度充足，且擅长处理中文的复杂语义。
  \item \href{https://chat.deepseek.com}{\textbf{DeepSeek}}：同样是优秀的国产工具，不仅注册简单无门槛，还在逻辑推理和代码生成上有亮眼表现，日常使用完全免费。
\end{itemize}

绝大多数平台的注册流程都非常类似，通常只需三步：

\begin{enumerate}
  \item 打开网站，点击「注册」或「Sign Up」。
  \item 填写邮箱或手机号，设置密码并完成验证。
  \item 登录后，你就能看到一个处于页面中央或底部的对话输入框。
\end{enumerate}

\begin{tipbox}
就如前面所说的，初学阶段不用纠结选哪个模型「最好」。工具只是载体，最重要的是掌握提示词的构建技巧。随意挑选一个顺眼的，现在就开始吧。
\end{tipbox}

\section{你的第一条提示词}

登录之后，你会看到一个输入框（通常在页面底部）。试试直接输入一句话：

\begin{promptbox}
帮我用三句话介绍杭州。
\end{promptbox}

按下回车（或发送按钮），等几秒钟，AI 就会给你一段回答。

你可能会得到类似这样的内容：

\begin{aibox}
杭州是浙江省的省会城市，以西湖闻名于世。这座城市有着两千多年的历史，是中国重要的电子商务中心。杭州也以龙井茶、丝绸和美食著称，是一座将古典韵味与现代活力完美融合的城市。
\end{aibox}

恭喜你，你已经完成了人生中的第一次 AI 对话。感觉是不是挺简单的（有手就会）？

不过你可能也会觉得：这个回答虽然没什么错误，但读起来像千篇一律的维基百科，没什么特别之处。
别着急，提示工程的魅力就在接下来的这一步。让我们试试换一种问法。

\section{同一个问题，换个问法试试}

现在，我们把同样的需求重新表达一遍，但这次加上具体的限制条件和细节要求：

\begin{promptbox}
你是一名旅游博主，请用轻松幽默的语气，用三句话向从没去过杭州的朋友介绍杭州，重点突出美食和自然风光。
\end{promptbox}

你可能会得到这样的回答：

\begin{aibox}
杭州是一座你去了就不想走的城市：西湖边散个步，感觉自己活成了古诗里的人。饿了就去河坊街，片儿川、叫花鸡、龙井虾仁轮番轰炸你的味蕾。如果你还没来过，那你的旅行清单上就少了最重要的一站。
\end{aibox}

对比一下两次回答，差别是不是很明显？

第一次的回答中规中矩，像百科全书；第二次活泼生动，像一个真正的旅游博主在跟你聊天。而我们做的唯一改变，就是在提示词里加了三样结构化约束条件：

\begin{enumerate}
  \item \textbf{角色设定（Role-playing）}：「你是一名旅游博主」
  \item \textbf{风格约束（Tone/Style）}：「轻松幽默的语气」
  \item \textbf{内容锚点（Focus/Content Anchors）}：「重点突出美食和自然风光」
\end{enumerate}

这就是提示词的「魔力」所在。当你没有给出任何约束时，模型会从它海量的训练数据中匹配最通用、最安全的回答方式。但当你给定了「角色」和「语气」等关键信息后，AI 就会把注意力集中在对应的风格和领域上。同一个 AI，能力没变，只是被你的提示词引导到了不同的方向，产出自然就不一样了。

\section{试试让 AI 帮你做一件真实的事}

理论讲完了，现在是你自己动手的时间。

请从你今天的待办清单里挑一件事，试着让 AI 帮你完成。比如：

\begin{itemize}
  \item 「帮我写一封给同事的邮件，通知下周三下午两点开项目复盘会议，语气正式但友好。」
  \item 「把下面这段英文翻译成中文，保持专业的语气：（粘贴你的英文内容）」
  \item 「帮我列一个周末两天的北京旅游计划，预算 2000 元以内，住宿和交通也帮我考虑进去。」
  \item 「我需要写一份季度工作总结，请帮我列一个大纲，包括工作成果、遇到的挑战和下季度计划。」
  \item 「我第一次请朋友来家里吃饭，帮我设计一桌 4 菜 1 汤的家常菜菜单，要有荤有素、好做又好看，买菜预算 100 元以内，并列出具体的食材采购清单。」
\end{itemize}

挑一个试试，真实地体验一下 AI 是如何帮你提高效率的。

如果它第一次给出的回答不够让你满意，完全没关系。你可以直接在对话框里追问和打磨，试着调整你的提示词：补充一些遗漏的细节、要求换一种说话方式，或者直接指令它「回答太冗长了，请精简并在 100 字以内总结要点」。

\begin{tipbox}
纸上得来终觉浅，绝知此事要躬行。亲自在对话框里敲下回车键，并观察 AI 随着你指令修改而发生的反馈变化，比阅读十页理论知识都要管用得多。
\end{tipbox}

% !TeX root = ../prompt-engineering-guide-for-beginners.tex

%%%%%%%%%%%%%%%%%%%%%%%%%%%%%%%%%%
%% 第三章：基本概念与模型设置
%%%%%%%%%%%%%%%%%%%%%%%%%%%%%%%%%%
\chapter{基本概念：理解大语言模型的工作方式}
\label{ch:basic-concepts}

在动手书写更复杂的提示词之前，我们先花几分钟了解一些大语言模型（LLM）的基本工作概念。
放心，这些概念并不艰涩，但掌握它们能帮助你理解 AI 底层的行为逻辑，从而写出更精炼、更有效的提示词。

\section{大语言模型的基本原理}

上一章我们提到，大语言模型本质上是一个「超级文字接龙选手」。这一节我们在此基础上稍微深入一点，介绍三个核心技术概念。

\subsection*{Token（令牌）}
\label{sec:token}

AI 不是像人类一样逐字逐句地阅读你输入的内容，而是先将文本切分成一个个更小的语义单元，这些单元在技术上叫做 Token。

对于英文来说，一个 Token 大约相当于一个完整的单词或单词的一部分（例如 "hamburger" 在某些模型引擎中会被 BPE（字节对编码）算法切分成 "ham" 和 "burger"）；对于中文来说，由于信息密度大，一个汉字通常对应 1 到 2 个 Token，具体取决于不同厂商底层模型训练时使用的词表（Vocabulary）切分与编码规则。

举个具体的例子来感受一下：「今天上海天气蛮好的」这句话大约会被切分成 7 到 12 个 Token（取决于使用的具体模型），而对应的英文 "The weather in Shanghai is great today" 大约需要 8 个 Token。各家模型厂商通常都提供了在线 Tokenizer 工具，你可以亲自输入一段文字，直观体验 Token 的切分过程。

\begin{tipbox}
为什么要深度了解 Token？因为在模型眼底，世界上没有“字”的概念，只有一长串对应的 Token ID 数字。模型对代码或长串数字的拆分往往违反人类直觉（例如把“3.1415”切成“3.14”和“15”），这就是产生弱智计算错误的根本原因。同时，所有大模型 API 接口计费、输入输出额度限制，全部都是以 Token 数量为标准计算的。
\end{tipbox}

\subsection*{上下文窗口（Context Window）}
\label{sec:context-window}

上下文窗口是 AI 在一次对话中能「看到」的信息总量上限，同样是用 Token 数量来衡量的。

你可以把它想象成为一张物理的「工作台」：工作台的面积是固定的，上面能摆放的文件也有限。只要不超过这个范围，AI 都能综合这些信息给出回答；但一旦你提供的信息超出范围，AI 就只能把相对「陈旧」的信息挤出工作台，结果就是它「忘记」了你之前的交代。

早期的模型上下文可能只有几千个 Token（约等于几页纸的内容），而目前的先进模型（例如 Kimi、Claude 3 等）有些已经支持几十万甚至上百万 Token（足以一次性塞入整整一本书的内容）。

\subsection*{补全（Completion）}
\label{sec:completion}

从技术底层来看，AI 进行信息生成的过程叫做「补全」。你在对话框内提交一段文字（即提示词），AI 所做的主要工作是根据统计概率，紧接着这段文字向下「续写」生成它认为最合理的后续内容。

即使你看起来是在以一问一答的方式发出问题，AI 其实也不在使用传统的思维去「思考」和「回答」，而是在计算这段对话接下来最可能出现怎样的后文。这就是为什么提示词起到了四两拨千斤的作用：你给的开头和背景哪怕只有微妙的区别，它续写演化的方向都会截然不同。

\section{大语言模型常规设置（LLM Settings）}

许多面向开发者和深度用户的 AI 工具（例如各种调用 API 的客户端、Playground 等）都会开放一些参数设置面板，允许你自定义或者干预模型的输出行为。掌握这些参数，你就能让 AI 变得按照你的特殊想法进行输出。

\subsection*{Temperature（温度）}
\label{sec:temperature}

Temperature 是出镜率最高的核心控制参数，它主导着大模型输出的「确定性与随机性」。数值一般从 0 开始。

\begin{itemize}
  \item \textbf{低温度值}（如 0.1 到 0.3）：剥夺 AI 的大部分发散能力。它的回答将变得非常确定、聚焦且极其稳定。该模式适合事实精准查询、结构化数据提取、代码生成、逻辑严密的数学计算等「只求准不求怪」的任务。
  \item \textbf{高温度值}（如 0.7 到 1.0 乃至更高）：解锁 AI 的天马行空。它的回答会充满创意与意想不到的表达方式，但相伴的风险就是更容易出现胡编乱造的「幻觉」。适合用在头脑风暴、虚构文学创作或是灵感启发的场景。
\end{itemize}

来感受一下 Temperature 对同一个问题的影响。假如你问 AI「用一句话形容上海」：

\begin{itemize}
  \item \textbf{低温度（0.2）}可能回答：「上海是中国最大的经济中心城市。」
  \item \textbf{高温度（0.9）}可能回答：「上海是一座让人在弄堂里吃着葱油拌面、却仰头就能看到陆家嘴天际线的魔法城市。」
\end{itemize}

同一个 AI，同一个问题，仅仅因为温度参数不同，输出的「画风」就判若两人。

\begin{tipbox}
口诀：追求靠谱选低温，想要创意推高温。
\end{tipbox}

\subsection*{Top-p（核采样）}

Top-p 是另一种常见的、用于调控输出概率多样性的参数。虽然它的技术原理与 Temperature 不同，但表现在结果上的作用极其相似。

Top-p 的设定值位于 0 到 1 之间。值越小，模型在选下一个词时就越拘谨（它只会从少数几个概率最高、最常出现的词汇中去进行选择）；值越大，那些原本冷门或平庸的词，也会有被选中的机会。

\begin{tipbox}
通常情况下，开发者们约定俗成的一条原则是：如果你想调整结果的丰富度，\textbf{只调整 Temperature 或只调整 Top-p 中的一个即可}，这两者没必要同时改动，否则容易让产出失控。
\end{tipbox}

\subsection*{Max Tokens（最大令牌数）}

此参数是一道红线，强制限定了 AI 单次回答时能生成的最大长度（Token 总量）。如果用户只想要一句简短的口号，把这个值压低，就能避免部分「过度热情」的大模型喋喋不休。

假设你将 Max Tokens 设为 100，当生成内容触及第 100 个 Token 时，不管句子是否讲完、标点符号有没有闭合，模型都会像被断掉电源一样立刻停止。这就要求你在实际应用中，不能一味图短而设得过于极致，以免最终得到被生硬截断的半成品回答。

\subsection*{Stop Sequences（停止序列）}

用户可以向系统中传入一个或多个特定的字符串列表（比如某种标记符号或某些句子），模型一旦在它的接龙过程中生成了与这个列表相匹配的内容，生成进程就会立即停下。

以多角色剧本生成为例，假如你希望每次生成一人份的台词，你可以将「李雷：」设定为停止序列。这样，系统在完成一段输出并准备开始「李雷：」时，就会被强行中断。如此一来，你将拿到单独一轮回合的干净文本，而不至于长篇大论。这个功能在代码提取或者 API 设计中常被作为截断垃圾文本的操作，虽然我们在日常文字聊天时鲜少触及。

\subsection*{Frequency / Presence Penalty（频率惩罚与存在惩罚）}

很多时候你可能遇到 AI 开始循环重复特定词语或陷入「车轱辘话」中，这两个参数就是应对该烦恼的克星。

Frequency Penalty（频率惩罚）的核心规则是按出现次数递减：模型在单次生成中每用过一次某词汇，该词再次出场的概率就在递减，重复越多打压越狠。
而 Presence Penalty（存在惩罚）奉行的则是「一视同仁的打压机制」：无论某个词汇出现了 1 次还是 10 次，只要「存在于本文中」，再次调用的门槛就会被施加一次固定的难度。

如果在实践中发现模型像口头禅一样频繁抛出「总而言之」、「另外一方面」（车轱辘话），你可以适当地调高这两个惩罚值，驱使模型调用更多新的词汇；代价则是如果调得过高，内容就会显得生硬和词不达意，因为它可能会为了避开常用词语而被「逼着」抛出生僻字词。

\subsection*{不同任务怎么调？}

下面这张速查表可以帮你快速决定参数设置：

\begin{center}
\begin{tabular}{lcc}
  \toprule
  \textbf{任务类型} & \textbf{Temperature} & \textbf{建议} \\
  \midrule
  事实问答、数据分析 & 0.1\textasciitilde 0.3 & 追求准确 \\
  邮件、报告、翻译    & 0.3\textasciitilde 0.5 & 稳定可靠 \\
  文案创作、头脑风暴  & 0.7\textasciitilde 0.9 & 鼓励创意 \\
  诗歌、小说、自由创作 & 0.9\textasciitilde 1.0 & 最大随机性 \\
  \bottomrule
\end{tabular}
\end{center}

如果你实在不知道怎么设置，保持默认值就好。大多数 AI 工具的默认参数已经针对通用场景做了优化。

\section{系统提示词（System Prompt）}
\label{sec:system-prompt}

许多现代大语言模型的应用（例如 ChatGPT、Claude 以及其他大语言模型生态软件如 Cherry Studio、AnythingLLM）都会专门提供一个功能：「系统提示词 / 系统预设」设定。

系统提示词的功能，其实是一道具有更高权限的背景环境预设墙。在它身上可以完成以下指令动作：

\begin{itemize}
  \item 奠定专业基础与人设：例如「你是一位有着 20 年经验的心胸外科医生，接下来的对话中要保持高度的同理心。」
  \item 控制核心风格：例如「始终只用严谨得体的规范中文输出，禁止使用网络缩写词及任何外语单词的混入使用。」
  \item 设置铁律与工作边界：例如「但凡遇到任何越界问题，你都必须立刻回复『很抱歉我不具备相关权限』，严禁任何擅自的脑补。」
\end{itemize}

你提交的这些系统预设，并不会夹杂在你平时的对话问答里。它隐匿于前端幕布之后，却是影响和重塑整个对话过程模型最终走向的最高统帅力量。

\begin{tipbox}
即使你使用的是没有专门设置入口的直接对话框，通过在第一次聊天发送的长提示词中直接带上一句「在接下来的整个会话窗口中，请将自己定位为一位富有经验的高级财务顾问」，这也是通过模拟系统提示词打下的工作基调。
\end{tipbox}

\section{多轮对话与上下文管理}
\label{sec:multi-turn}

绝大多数场景中，大家感受不到 AI 在和自己互动沟通中产生断层的违和感，因为它能承接上文的语境。这是由于 AI 在幕后做了一件体力活，它把从对话开启至今你发出的每一条记录和它做出的每一句回答全部放在了自己的上下文窗口中。

遗憾的是，这种记忆方式绝非无限容量。正如前面小节解释的，我们已经知道了它是把记忆放置在「一定尺寸的物理工作台」上的。如果你们在这间聊天室中不断对答，产生的内容在字数上限度超过了当前平台所规定的窗口极限尺寸时，系统便会自动启动机制，将最先产生的对话（通常指的是第一句话和后续的早期互动内容）强行抹去，来为新生内容让路。这就造成了我们在连续对话数日后发现，AI “失去记忆”，完全忘记了当初角色扮演设定。

面对这一大模型的底层短板，目前有以下几点有效的提示工程应对小窍门：

\begin{itemize}
  \item \textbf{尽早开启新对话分支（Reset Session）}：如果你目前的交流与先前的分析毫无关系，请务必「另起炉灶」清空当前记录（重启 Chat Session）。这样不但可以还 AI 一块空白的“无污染算力画板”，更能避免由于上下文过于庞大导致 AI 注意力分散和 API 额度浪费。
  \item \textbf{利用阶段性回溯进行状态压缩（State Compression）}：如果是跨越长文本的写作或分析过程。在感到上下文快吃满前，请让 AI 执行：「基于以上沟通进展，请输出一份结构化总结」。随后你可以开启新对话，并把这份“高纯度压缩包”作为新环境的起始状态（Prompt），直接跨过冗长无用的聊天记录。
  \item \textbf{重申关键铁律信息}：面对开始偏离轨道的 AI，与其在对话框内去用自然语言与它反复争辩、指责它忘记旧规则（这反而会在上下文中注入更多污染性的 Token），不如直接在新的命令前再叠加一次最高权重的提示内容：「请再次使用开局定义的规则去作回复」。
  \item \textbf{对抗“迷失在中间（Lost in the Middle）”效应}：多项论文研究表明，当上下文达到数万 Token 时，模型抓取顶端和尾部信息的注意力分数始终保持最高，而中间段落的信息极易被忽略（即所谓的 U 型注意力曲线）。因此如果有影响存亡的规矩、或者极为核心的输出格式参考，请务必将其排布在整个 Prompt 的最顶部（作为 System 预设）或者最末尾的总结请求处。
\end{itemize}

系统参数、基础工作概念和记忆法则这三块基础知识，是你后续运用各种提示技巧时最重要的底座支撑。

\section{补充：推理模型与「深度思考」}
\label{sec:reasoning-models}

近一两年，你可能听说过一类特别的模型，如 OpenAI 的 o1/o3 系列、DeepSeek R1 等。它们被称为「推理模型」或「深度思考模型」，和普通的对话模型有一个关键区别。

普通对话模型在收到你的提示词后，会直接开始逐个生成回答的每一个 Token。而推理模型会在生成最终回答之前，先在内部进行一段较长的「思考过程」（Think / Chain of Thought），把问题拆解、验证、反复推敲之后，再给出一个更加深思熟虑的答案。

用一个直观的类比来理解：普通对话模型像是一个脱口而出的即兴演讲者，话到嘴边就说出来；而推理模型更像是一位考场上的考生，先在草稿纸上列出思路、验算一遍，确认无误后才把答案誊写到答题卡上。

\textbf{来看一个具体的例子。}假设你问 AI 这样一道逻辑题：

\begin{promptbox}
一个房间里有 3 个开关，分别控制隔壁房间的 3 盏灯。你只能进入隔壁房间一次。请问如何确定每个开关分别对应哪盏灯？
\end{promptbox}

如果使用普通对话模型，它可能会给出一个笼统的回答，甚至直接跳到结论而漏掉关键的推理步骤。但推理模型会在内部展开一段完整的思考链（你通常可以在界面上看到折叠的「思考过程」）：先分析约束条件（只能进入一次），再逐步推导出利用灯泡温度作为额外信息的策略，最终给出严密的三步解法。

\textbf{再看一个数学场景。}你让 AI 计算：

\begin{promptbox}
一个水池有两个进水管和一个出水管。A 管单独注满需要 6 小时，B 管单独注满需要 8 小时，出水管单独放完需要 12 小时。三管同时打开，多久能注满水池？
\end{promptbox}

普通模型可能直接给出一个数字却在中间步骤中犯运算错误。推理模型则会将问题拆解为分数运算，列出 $\frac{1}{6} + \frac{1}{8} - \frac{1}{12}$ 的计算过程，通分、化简，最终得出正确答案，并且每一步都可以被验证。

\textbf{什么时候用推理模型？}

\begin{itemize}
  \item 数学计算、逻辑推理等需要严密推导的任务
  \item 复杂的代码编写和调试（例如涉及多个模块的架构设计或难以定位的 Bug 排查）
  \item 需要多步骤分析和决策的问题（例如商业方案的利弊权衡、多约束条件下的方案选择）
  \item 需要对自身回答进行自我检验和纠错的场景
\end{itemize}

\textbf{什么时候用普通模型就够了？}

\begin{itemize}
  \item 日常对话、文案写作、翻译
  \item 简单的问答和信息整理
  \item 对响应速度要求较高的场景（推理模型由于需要经历「思考」阶段，通常响应更慢，消耗的 Token 也更多，费用相应更高）
\end{itemize}

\begin{tipbox}
使用推理模型时，提示词反而不需要太复杂。你不必像普通模型那样加上「请一步一步思考」之类的引导语，因为它本身就会自动进行深度推理。只需要把问题描述清楚、把约束条件列全即可。过多的格式要求或引导性话术反而可能干扰模型自身的推理节奏。
\end{tipbox}

理解了这些基础概念之后，接下来我们将开始学习如何组装提示词的各个要素。

% !TeX root = ../prompt-engineering-guide-for-beginners.tex

%%%%%%%%%%%%%%%%%%%%%%%%%%%%%%%%%%
%% 第四章：提示词要素
%%%%%%%%%%%%%%%%%%%%%%%%%%%%%%%%%%
\chapter{提示词要素：一条好 Prompt 的解剖课}
\label{ch:prompt-elements}

了解了 AI 的基本工作原理之后，我们来聊聊提示词到底应该「长什么样」。
一条好的提示词并不是随手一写那么简单，它有自己内在的结构和要素。
搞清楚这些要素，你就掌握了写好提示词的底层方法论。

\section{提示词的四大要素}

一条完整的提示词，最多可包含四个核心组成部分。虽然不是每次对话都要全部写满，但了解它们能帮你在需要时，像搭积木一样快速组装出高质量的指令。

\subsection*{指令（Instruction）}

这是提示词的核心：明确告诉 AI 你想让它做什么。

好的指令要用具体的动词，避免笼统含糊。比如：

\begin{itemize}
  \item 「总结」而不是「说说」
  \item 「列出 5 个要点」而不是「随便讲讲」
  \item 「翻译成英文」而不是「帮我弄一下」
\end{itemize}

\subsection*{上下文（Context）}

背景信息能够帮助 AI 更好地理解你的需求。你提供的上下文越充分，AI 的回答就越贴合你的实际应用场景。

比如同样是「帮我写一封邮件」，在这基础上加上「这是发给甲方的催款邮件，我们合作了三年，关系一直不错，但这次对方拖了两个月没付款」，AI 就能极其精准地拿捏住「既不撕破脸又要施加压力」的妥帖语气。

\subsection*{输入数据（Input Data / Payload）}

如果你需要 AI 处理一段具体的内容，就需要把这段内容干净地注入提示词中。比如：

\begin{itemize}
  \item 需要处理提取知识点的数百行源文件代码
  \item 需要总结的文章文本（Context Corpus）
  \item 需要做可视化推理的脏数据表（CSV）
  \item 需要改写的初稿段落
\end{itemize}

\begin{tipbox}
从工程化角度来看，Prompt 就是一段由“操作代码 (Instruction)”和“数据参数 (Payload)”组成的代码块。为了防止 AI 在读取时把你的文章内容误认为是新下达的指令（即常见的 Prompt Injection 提示词注入漏洞），强烈建议使用特定的 XML 标签（如 \verb|<text>...</text>|）或三个反引号（\verb|```|）将输入数据严格包裹隔离起来。
\end{tipbox}

\subsection*{输出格式（Output Format / Schema Constraints）}

告诉 AI 你期望回答以何种结构形式呈现。作为严谨的提示词工程师，你应当将其理解为与 API 协同中的接口 Schema 定义：

\begin{itemize}
  \item 「以 Markdown 表格形式展示」
  \item 「用层级清晰的编号列表列出」
  \item 「仅仅输出合法的 JSON 格式字符串，不需要任何解释」
  \item 「使用 LaTeX 公式语法包裹所有数学计算过程」
  \item 「强硬控制输出不超过 300 个 Token」
\end{itemize}

\section{万能公式：角色 + 背景 + 任务 + 格式要求}
\label{sec:prompt-formula}

如果你觉得四大要素不太好记，这里有一个更简单的万能公式。写提示词的时候，按照这个顺序填空就行：

\begin{enumerate}
  \item \textbf{角色}：让 AI 扮演什么身份？
  \item \textbf{背景}：这件事的背景是什么？
  \item \textbf{任务}：你需要 AI 做什么？
  \item \textbf{格式要求}：你希望回答是什么样子的？
\end{enumerate}

来看一个实际例子。

\textbf{普通版：}

\begin{promptbox}
帮我写个请假条。
\end{promptbox}

AI 可能会写出一个非常模板化、没有细节的请假条。

\textbf{万能公式版：}

\begin{promptbox}
你是一名 500 强企业的员工（角色），由于重感冒需要向严厉的老板请假两天（背景），请帮我写一封请假邮件（任务），要求态度诚恳、措辞专业、字数 100 字以内（格式要求）。
\end{promptbox}

用了万能公式之后，AI 会给出一封语气恰当、内容具体、长度合适的请假邮件。

不需要每次都写得这么面面俱到，但当你对 AI 的回答感到不满时，可以回过头来拿着这个公式检查一下：是不是缺了角色？是不是没交代背景？还是没说清楚输出的格式？

再来一个接地气的例子：

\begin{promptbox}
你是一位上海本地美食达人（角色），我外地朋友下周来上海玩三天（背景），请推荐五家不踩雷的本帮菜馆，涵盖从人均 50 到人均 300 的不同档次（任务），用表格列出餐厅名、地址、人均消费和推荐菜品（格式要求）。
\end{promptbox}

这样一写，AI 就不会给你推荐一堆连锁快餐或者搜索结果里千篇一律的「十大必吃」了，而是会按照你要求的表格格式、价位梯度，给出更有参考价值的推荐。

\section{好提示词 vs 坏提示词：对比拆解}

下面通过几组对比，让你直观感受提示词质量对结果的影响。

\subsection*{案例一：写一篇产品介绍}

\textbf{坏提示词：}
\begin{promptbox}
帮我写个产品介绍。
\end{promptbox}

\textbf{问题}：AI 不知道是什么产品，也不知道写给谁看，只能写出一篇空洞的通用模板。

\textbf{好提示词：}
\begin{promptbox}
请帮我写一段产品介绍。产品是一款面向年轻白领的便携式咖啡机，主打卖点是 30 秒快速萃取和超轻便携设计。介绍文案发布在小红书上，语气活泼种草，控制在 150 字以内。
\end{promptbox}

\textbf{差别}：有了具体产品、目标人群、发布渠道、语气和字数要求，AI 就能精准输出。

\subsection*{案例二：翻译一段文字}

\textbf{坏提示词：}
\begin{promptbox}
翻译这个：We need to leverage our core competencies to drive synergistic outcomes.
\end{promptbox}

\textbf{好提示词：}
\begin{promptbox}
请把下面这句英文翻译成中文，保持商务风格，避免直译，用自然流畅的表达：

We need to leverage our core competencies to drive synergistic outcomes.
\end{promptbox}

\textbf{差别}：加上翻译风格和要求，AI 的翻译会更加自然地道，而不是逐字硬译。

\subsection*{案例三：做一个工作计划}

\textbf{坏提示词：}
\begin{promptbox}
帮我做个计划。
\end{promptbox}

\textbf{好提示词：}
\begin{promptbox}
我是一名新媒体运营，下个月需要完成公司公众号的日常运营。请帮我制定一份为期 4 周的内容发布计划，要求每周发布 3 篇文章，涵盖行业资讯、产品介绍和用户故事三个主题，以表格形式呈现，包含日期、主题、标题建议和关键词。
\end{promptbox}

\textbf{差别}：拥有了详细的背景交代、清晰的任务目标和明确的表格输出格式，AI 便能直接产出一份可落地执行的高质量计划。

通过这几组对比可以看出，写好提示词的核心其实就是三个字：\textbf{说清楚}。你给 AI 的信息越清晰、越具体，它的输出就越贴近你的预期。

\section{快速检查清单}

每次写完提示词、准备发送之前，可以用下面这个清单快速自查一遍：

\begin{center}
\begin{tabular}{lp{9cm}}
  \toprule
  \textbf{检查项} & \textbf{问自己} \\
  \midrule
  任务明确 & 我是否用了具体的动词（总结/翻译/列出），而不是「帮我看看」？ \\
  背景充分 & AI 是否知道这件事的前因后果和使用场景？ \\
  受众清晰 & AI 是否知道这个输出是给谁看的？ \\
  格式要求 & 我是否说明了输出的形式（表格/列表/字数）？ \\
  约束到位 & 我是否给出了语气、长度、风格等限制条件？ \\
  \bottomrule
\end{tabular}
\end{center}

不必每次都严格检查所有项目，简单的问题一两句话就够了。但当 AI 的回答不尽如人意时，回来对照这个清单，通常很快就能找到原因。

现在你已经知道了一条好提示词应该包含哪些「结构要素」，接下来我们将进一步聊聊该「如何表达」。下一章会为你介绍一系列实用的通用写作技巧。

% !TeX root = ../prompt-engineering-guide-for-beginners.tex

%%%%%%%%%%%%%%%%%%%%%%%%%%%%%%%%%%
%% 第五章：设计提示的通用技巧
%%%%%%%%%%%%%%%%%%%%%%%%%%%%%%%%%%
\chapter{设计提示的通用技巧：写出高分 Prompt 的方法论}
\label{ch:general-tips}

上一章我们拆解了提示词的结构，现在来聊一些经过大量实践验证的实用技巧。
这些技巧不针对某个特定场景，而是适用于所有类型的提示词。
掌握了这些方法论，你写提示词的水平会有一个明显的跃升。

\section{动词要具体，拒绝假大空}

很多人写提示词的习惯是用「说一下」「帮我看看」「讲一讲」这种笼统的动词。AI 拿到这样的指令，往往不知道你想要什么样的输出，只能给一个泛泛的回答。

解决办法很简单：换成更精确的动词。

\begin{center}
\begin{tabular}{ll}
  \toprule
  \textbf{模糊动词} & \textbf{精确动词} \\
  \midrule
  说一下 & 列出、总结、解释、分析 \\
  帮我看看 & 检查、对比、评估、纠错 \\
  讲一讲 & 描述、定义、举例说明 \\
  弄一下 & 改写、翻译、格式化、提炼 \\
  \bottomrule
\end{tabular}
\end{center}

例如：

\begin{itemize}
  \item 「说说这个产品的优缺点」→「\textbf{列出}这个产品的 3 个优点和 3 个缺点，每个配一句简短的解释」
  \item 「帮我看看这段代码」→「\textbf{检查}这段代码是否存在性能问题，并给出优化建议」
\end{itemize}

一个更精确、聚焦的动词，就能让 AI 瞬间领会侧重点，从而将回答质量提升一整个档次。

\section{提供上下文：AI 不会读心术}

AI 只能根据你告诉它的信息来回答。你脑子里知道的背景，如果没有写进提示词，AI 就是不知道。

所以，写提示词的时候要养成一个习惯：把相关的背景信息尽量提供给 AI。

\begin{itemize}
  \item 你是谁（角色/身份）
  \item 你的受众是谁（给谁看/用的）
  \item 你的目标是什么（想达到什么效果）
  \item 有什么限制条件（公司规定、预算、时间等）
  \item 有没有参考资料（是的话，直接贴进来）
\end{itemize}

\textbf{反面例子：}
\begin{promptbox}
帮我写一段欢迎词。
\end{promptbox}

\textbf{正面例子：}
\begin{promptbox}
我是一家少儿编程培训机构的校长，下周六有一场面向 6 到 12 岁小朋友家长的开放日活动。请帮我写一段 3 分钟的欢迎词，语气热情但不夸张，重点介绍我们的教学特色和孩子们能收获什么。
\end{promptbox}

信息越是生动丰满，AI 的回答就越像是一个深入一线、懂你业务的资深同行执笔的，而不是一个走过场敷衍了事的机器路人。

\section{设定约束条件}

不给约束条件的提示词就像是一张没有边界的白纸，AI 不知道该往哪个方向写，只能写一大堆「正确但没用」的内容。

常见的约束条件包括：

\begin{itemize}
  \item \textbf{字数限制}：「控制在 200 字以内」「不超过 3 段」
  \item \textbf{语气要求}：「正式/轻松/幽默/严肃」
  \item \textbf{受众定义}：「写给完全没有技术背景的人看」
  \item \textbf{排除条件}：「不要使用专业术语」「不要列举超过 5 个」
  \item \textbf{格式限定}：「以 Markdown 表格形式输出」「用编号列表」
\end{itemize}

约束条件越是详细，AI 的产物就越贴合你的心意。你可以把约束条件想象成预先给 AI 画好了一个边界分明的工作框：在这个框以内，它拥有极大地文字组织自由；但框外那些「正确但你不需要」的废话，统统被强制过滤掉了。

\section{分步拆解复杂任务}

如果面对的需求本身异常庞大复杂，试图用“一言以蔽之”的方式把所有指令全塞进对话框里是非常危险的，往往会导致 AI 在生成时顾此失彼。
针对此类场景，最稳妥的策略就是化繁为简：把一个超级任务肢解为先后承接的几个小短步。

\textbf{不太好的做法：}
\begin{promptbox}
帮我写一篇关于远程办公趋势的深度文章，要有数据引用、案例分析、未来展望，3000 字以上。
\end{promptbox}

\textbf{更好的做法：}

\begin{promptbox}
第一步：请帮我列出关于「远程办公趋势」的文章大纲，包括 5 到 7 个主要章节。

第二步（拿到大纲并确认没问题后）：请基于以上大纲的指引，逐章推进撰写。请先集中精力只完成第一章的内容：远程办公的现状与底层数据支撑。

第三步（如此循环，直到全文章节完成）：现全文已生成完毕，请你作为最终主编通读全文，集中修正可能存在的逻辑断层，并专门润色章节之间的过渡段，使全文阅读如同行云流水。
\end{promptbox}

分步拆解最大的优势在于：每个独立小步骤的任务边界都极其明确清晰，AI 在局部的正确率和表现力上就会达到峰值。更重要的是，在每一轮微小动作交付后，你都能及时进行干预修正，而不必面临“等了 3 千字出来却发现必须全盘推翻重写”的惨痛沉没成本。

\section{使用分隔符组织内容}

当你的提示词里既有指令又有输入数据时，用分隔符把它们清晰地隔开，AI 就不会搞混。

常用的分隔符有：

\begin{itemize}
  \item 三个引号：\verb|"""|
  \item 三个反引号：\verb|```|
  \item XML 标签：\verb|<text>...</text>|
  \item 横线：\verb|---|
\end{itemize}

\textbf{不用分隔符的例子：}
\begin{promptbox}
请总结以下内容的核心观点远程办公正逐渐成为许多企业的常态化选择根据最新调查显示超过60\%的员工希望每周至少有两天在家办公……
\end{promptbox}

AI 可能分不清哪些是你的指令、哪些是需要处理的内容。

\textbf{用了分隔符的例子：}
\begin{promptbox}
请总结以下内容的核心观点：

\verb|"""|

远程办公正逐渐成为许多企业的常态化选择。根据最新调查显示，超过 60\% 的员工希望每周至少有两天在家办公……

\verb|"""|
\end{promptbox}

用分隔符把指令和数据区分开后，AI 就能清楚地知道哪些是你的操作指令、哪些是需要处理的原始素材，从而避免混淆。

\section{先说「做什么」而非「不做什么」}

人在提需求时，常常会习惯性地使用否定句式来做限制，比如「不要写得太长」「不要使用专业术语」「不要太正式」。

这里面有个有趣的现象：当你告诉 AI「不要写得太长」时，模型内部反而会先激活与「太长」相关的概念，这可能导致它偏向那些你想要避免的方向。相比之下，直接用正向指令（明确说出你想要什么）的效果通常更好。道理很简单：你告诉了 AI 不要做什么，但它仍然不知道你到底想要什么。

\begin{center}
\begin{tabular}{ll}
  \toprule
  \textbf{负向指令} & \textbf{正向指令} \\
  \midrule
  不要写得太长 & 控制在 150 字以内 \\
  不要用专业术语 & 用通俗易懂的语言 \\
  不要太正式 & 用轻松的口语化语气 \\
  不要列举太多 & 只列出最重要的 3 点 \\
  \bottomrule
\end{tabular}
\end{center}

当然，绝对性禁令在某些底线场合依旧是无可替代的（例如系统指令里规定的「绝对禁止编造虚假数据并以此对用户作答」）。但是在 90\% 的日常指令沟通下，试着切换心态，先把「鼓励它走向什么模样」描绘得清清楚楚，成效必然会截然不同。如果它仍然有跑偏的迹象，你再在句尾补上一刀轻微的反向封禁即可。

到这里，你已经掌握了写好提示词的通用方法论。接下来的三章，我们将按照「入门 $\to$ 进阶 $\to$ 高阶」的梯度，为你介绍提示工程领域中几种经过验证的具体技巧。别被名字吓到，其中很多方法你可能已经在实践中无意识地用过了。

%!%
% Copyright (c) 2026 mingcheng <mingcheng@outlook.com>
% 
% File: ch06-prompting-basics.tex
% Author: mingcheng <mingcheng@outlook.com>
% File Created: 2026-03-27 15:50:02
% 
% Modified By: mingcheng <mingcheng@outlook.com>
% Last Modified: 2026-03-29 07:50:18
%%

% !TeX root = ../prompt-engineering-guide-for-beginners.tex

\chapter{提示工程入门：最常用的四把钥匙}
\label{ch:prompting-basics}

前面几章我们打好了基础，现在正式进入提示工程的技巧部分。
这一章介绍的四种技巧是日常使用频率最高的，掌握它们就能覆盖你 80\% 以上的使用场景。

\section{零样本提示（Zero-Shot Prompting）}
\label{sec:zero-shot}

这是最简单直接的提示方式：不给任何示例，直接告诉 AI 你的需求。

\textit{类比：你跟一个聪明的实习生说「帮我翻译这段话」，他直接就能做，不需要你先示范。}

在机器学习领域，由于大语言模型在预训练阶段经过了海量任务数据的微调（Zero-Shot Transfer），它们已经获得了强大的“任务泛化（Generalization）”能力。所以对于很多人类常识范围内的任务，你不需要额外喂给它案例集，直接下达 System 级指令就好。

\textbf{示例：}

\begin{promptbox}
请把下面这段文字翻译成英文：

今天天气真好，适合出去走走。
\end{promptbox}

\begin{promptbox}
请将以下文字的情感分类为「正面」「负面」或「中性」：

这家餐厅的菜品味道一般，但服务态度还不错。
\end{promptbox}

零样本提示适用于\textbf{任务需求边界极度清晰、且无需在最终表现形式上做任何包装点缀的强功能性场景}。比如：单纯的中外文互译、常识科普型问答、文本的基础情绪归类、或者干脆把 AI 当成高级语病检查器来使用。

但很快你就会在实践中发现：当你期待 AI 的输出不仅要完成任务，还要带有特定的排版和风格（例如你公司特有的月报格式、或者某种特定的文风）时，零样本提示就不太够用了。这时候就需要用到下一种技巧：少样本提示。

\section{少样本提示（Few-Shot Prompting / In-Context Learning）}
\label{sec:few-shot}

有时候光靠语言描述很难讲清楚你到底想要什么样的输出。这时候最有效的办法是：直接给 AI 看几个例子，触发大模型的“上下文学习（In-Context Learning）”能力。

\textit{类比：你给实习生看了三份合格的周报样本，他立马知道格式、语气和重点该怎么写。}

\textbf{示例：}

\begin{promptbox}
请根据以下示例，为新的产品生成一句广告语：

产品：运动手环

广告语：24 小时守护你的每一次心跳

产品：降噪耳机

广告语：喧嚣世界里，给你一片宁静

产品：智能台灯

广告语：
\end{promptbox}

AI 会根据前面两个例子的模式，为「智能台灯」生成风格一致的广告语。

少样本提示的几个要点：

\begin{itemize}
  \item \textbf{样本数量}：通常提供 2 到 5 个示例即可。太少 AI 可能抓不住规律；太多则浪费上下文窗口的空间，甚至可能导致 AI 过度模仿而失去灵活性。
  \item \textbf{样本质量}：你提供的示例就是 AI 模仿的「标准答案」，所以务必确保每个示例本身在内容和表达上都是你满意的状态。
  \item \textbf{格式统一}：所有示例的结构、标点符号的使用方式、段落排版都应保持高度一致，像模板一样整齐。
  \item \textbf{覆盖多样性}：几个示例之间尽量涵盖不同的情况和类型，这样能帮助 AI 更好地理解规律并举一反三。
\end{itemize}

少样本提示特别适合以下场景：自定义格式的输出、特定风格的写作、分类任务需要定义类别的情况、以及你难以用语言精确描述期望输出的时候。

\section{链式思考提示（Chain-of-Thought, CoT）}
\label{sec:cot}

当你让 AI 处理需要推理的任务（比如数学题、逻辑判断、多步骤分析），直接问它要答案，它可能会给出错误的结果。但如果你让它「一步一步地想」，它的准确率会大幅提升。

\textit{类比：让实习生别直接报答案，先在纸上写出解题步骤。中间出错你也能看到并纠正。}

\textbf{不用链式思考的例子：}

\begin{promptbox}
小明有 15 个苹果，他给了小红 3 个，又从小华那里拿了 7 个。后来他把手里苹果的一半分给了妈妈。小明现在有几个苹果？
\end{promptbox}

AI 可能直接给一个答案，但不一定对。

\textbf{用了链式思考的例子：}

\begin{promptbox}
小明有 15 个苹果，他给了小红 3 个，又从小华那里拿了 7 个。后来他把手里苹果的一半分给了妈妈。小明现在有几个苹果？

请一步一步思考，展示你的推理过程。
\end{promptbox}

AI 会这样回答：
\begin{aibox}
第一步：小明开始有 15 个苹果。

第二步：给了小红 3 个，剩下 15 - 3 = 12 个。

第三步：从小华那里拿了 7 个，现在有 12 + 7 = 19 个。

第四步：把一半分给妈妈，19 / 2 = 9.5，取整为 9 个（如果不能分半个苹果的话）。

所以小明现在有 9 个半苹果（或者 9 个，取决于苹果能不能切开）。
\end{aibox}

使用链式思考的好处不仅是提高最终答案的准确率。因为模型展示了完整的推理过程，一旦某一步出现了逻辑错误，你就能迅速定位问题所在，然后在对话中要求它从出错的那一步重新推导。

使用链式思考提示的触发咒语极其简练，你甚至只需要在那些需要长距离逻辑折返的问题行末尾打上一句简单的：

\begin{itemize}
  \item 「请一步一步仔细思考」（泛用度极高的经典咏唱）
  \item 「请向我展示得出该推论的核心中间验证过程」
  \item 「请先对材料进行逐条剖析后再给出汇总结论」
  \item 「Let's think step by step」（玄学上的大招，这句英文在早先大基座模型实验效果上的提分异常显著，不妨一试）
\end{itemize}

链式思考提示适合所有需要逻辑推理的场景：数学计算、因果分析、利弊权衡、决策建议等。

\section{角色扮演提示（Role Playing）}
\label{sec:role-playing}

给 AI 指定一个角色身份，它的输出风格、专业度和关注点都会跟着变化。这个技巧简单但非常好用。

\textit{类比：告诉实习生「你现在假装是客户来投诉」，他的语气和关注点立刻不一样。}

\textbf{示例：}

\begin{promptbox}
你是一位拥有 20 年经验的资深营养师。我最近经常熬夜加班，饮食不规律，请给我一份适合上班族的一周健康饮食建议。
\end{promptbox}

\begin{promptbox}
你是一名专业的产品经理，擅长用户需求分析。请帮我分析以下用户反馈，提炼出三个最核心的需求：

（粘贴用户反馈内容）
\end{promptbox}

\begin{promptbox}
你是一位严厉但公正的面试官，请针对「Java 后端开发工程师」这个岗位，给我出 5 道技术面试题，并附上参考答案。
\end{promptbox}

角色扮演的一些技巧：

\begin{itemize}
  \item \textbf{角色越具体越好}：「资深营养师」比「营养方面的专家」更有效，「拥有 20 年经验的儿科医生」比「医生」更有效
  \item \textbf{可以叠加多个角色特征}：「你是一名擅长通俗化表达的科学家」，这样既有专业度又有可读性
  \item \textbf{可以让 AI 扮演你的对立面}：比如让它扮演「挑刺的客户」或「试图找漏洞的黑客」来帮你检查方案的薄弱环节
  \item \textbf{有趣的角色能激发创意}：让 AI 扮演「一个 5 岁小孩」来解释复杂概念，或扮演「毒舌评论家」来吐槽你的文案
\end{itemize}

这四种技巧（零样本、少样本、链式思考、角色扮演）是提示工程中使用频率最高的基础工具。在实际工作中，你完全可以把它们灵活组合使用。比如：

\begin{promptbox}
你现在是一位在华尔街工作多年的对冲基金经理（角色扮演），请一步一步分析以下几组备选方案的潜在风险（链式思考）。
\end{promptbox}

掌握了这四种基础技巧，你已经可以应对 90\% 以上的日常办公、文本处理和信息获取场景了。

\begin{tipbox}
快速选择指南：如果任务简单直接，先试零样本；如果 AI 理解不到位，加几个示例走少样本；如果需要严密推理，加上「请一步一步思考」开启链式思考；如果想要特定风格或专业视角，给 AI 分配一个角色。这四把钥匙可以单独使用，也可以自由组合。
\end{tipbox}

如果你还想让 AI 处理更复杂的任务，比如需要极高准确性的推理、多步骤协作等，下一章会介绍几种更进阶的增强技巧。

% !TeX root = ../prompt-engineering-guide-for-beginners.tex

%%%%%%%%%%%%%%%%%%%%%%%%%%%%%%%%%%
%% 第七章：提示工程进阶篇
%%%%%%%%%%%%%%%%%%%%%%%%%%%%%%%%%%
\chapter{提示工程进阶：让回答更靠谱的增强技巧}
\label{ch:prompting-intermediate}

上一章介绍的基础四板斧，已经能应对你在办公和生活中的多数日常任务。
但如果你遇到需要极其严谨或关乎重要决策的任务，并且想让 AI 的回答彻底摆脱偶尔的“胡说八道”，就需要用到下面这六种进阶增强技巧。
它们是基础技巧的「升级补丁」，能大幅提升模型的可靠性。

\section{自我一致性（Self-Consistency）}
\label{sec:self-consistency}

上一章我们学到了链式思考（CoT，参见第~\ref{sec:cot}~节），通过让 AI 展示推理过程来提高准确率。但即便如此，一条孤立的思维链依然有可能「想偏」，一旦前面的逻辑出现谬误，后面的推理就会全线崩溃。

自我一致性就是用来对冲这层风险的：它的做法是让 AI 针对同一个问题，在底层通过调整 Temperature（温度树）或者 Top-p 参数，强行走出多条不同的推导计算分支（解码路径），然后利用集成学习式的“多数投票（Majority Vote）”机制，对得出的多个最终大纲进行比对，采纳绝大多数路径都指向的那一个答案。

\textit{类比：同一道数学题让实习生用三种方法做，三种方法都得到同一个答案，你就更放心了。}

\textbf{使用方法：}

\begin{promptbox}
以下是一个逻辑推理题：

一家公司有 100 名员工，其中 60\% 会骑自行车，70\% 会游泳，80\% 会打篮球。请问至少有多少员工同时会这三项技能？

请用三种不同的推理方法来解答这道题，然后对比三种方法的结果，给出最终答案。
\end{promptbox}

当你面对重要决策或复杂推理题时，使用自我一致性可以显著降低 AI 出错的概率。

\section{生成知识提示（Generated Knowledge Prompting）}

面对相对生僻或专业度极高的问题，直接发问很容易迫使 AI 开始“盲答”或捏造。此时，更理智的策略是：先命令 AI 去它庞大的底层数据库中，将相关的专业知识点系统性地提取出来，完成一次「闭卷默写」，然后再让它拿着这份自己默写出来的知识库去推演解答。

\textit{类比：考试前让实习生先把知识点默写一遍，再开始答题，正确率更高。}

这个技巧分两步：

\textbf{第一步：生成知识}
\begin{promptbox}
请列出关于「地中海饮食」的 5 个核心知识点，包括其定义、主要食物组成、健康益处、适用人群和常见误区。
\end{promptbox}

\textbf{第二步：基于知识回答问题}
\begin{promptbox}
基于以上知识点，请帮一位 45 岁、有轻度高血压的上班族制定一份为期一周的地中海饮食方案。
\end{promptbox}

通过两步走，AI 的回答会更加有理有据，而不是凭空编造。

当然，你也可以把两步合并到一条提示词里：

\begin{promptbox}
请先回顾关于「地中海饮食」的核心知识，然后基于这些知识，为一位 45 岁有轻度高血压的上班族制定一周饮食方案。
\end{promptbox}

\section{提示链（Prompt Chaining）}

有些任务太复杂，一条提示词搞不定。提示链的思路是：把一个大任务拆成多个小步骤，每一步是一条独立的提示词，前一步的输出直接作为下一步的输入。

\textit{类比：不是让实习生一口气写完一份商业计划书，而是先做市场调研，再写竞品分析，然后定策略，最后组装成文。}

\textbf{实际案例：写一篇产品评测文章}

\begin{promptbox}
\textbf{第一步}：「请帮我列出 iPhone 16 的 10 个核心特性，按照硬件、软件、设计三个维度分类。」

\textbf{第二步}：「基于上面的特性列表，请从一个普通消费者的角度，对每个特性给出优点和不足的评价。」

\textbf{第三步}：「请把上面的分析整合成一篇 1500 字的产品评测文章，结构为：开头引言、核心亮点、值得吐槽的地方、购买建议、总结。语气客观但不无聊。」
\end{promptbox}

提示链的优势在于：

\begin{itemize}
  \item 每一步的任务都足够简单，AI 更容易做好
  \item 你可以在每一步结束后检查结果，发现问题及时纠正
  \item 整体输出的质量通常比一步到位要高很多
\end{itemize}

\section{方向性刺激提示（Directional Stimulus Prompting）}

有时候你手头有一份繁杂的初始材料，你既不想直接把答案喂到 AI 嘴里限制它的发挥，又担心它的思路宛如脱缰野马。这时候，方向性刺激提示就能派上用场：在提示词的结尾埋入一个「高亮线索」或「思考引子」，牵引 AI 的视觉重心。

\textit{类比：侦探导师不直接告诉学徒凶手是谁，而是指了指带血的门把手说「去查查这个把手的指纹」，学徒就能顺藤摸瓜。}

\textbf{不加暗示：}
\begin{promptbox}
请分析这家公司最近业绩下滑的原因。

（附上公司的财务报告）
\end{promptbox}

\textbf{加了暗示：}
\begin{promptbox}
请分析这家公司最近业绩下滑的原因。重点关注研发支出的变化趋势和新品上市节奏。

（附上公司的财务报告）
\end{promptbox}

通过给出一些方向性的提示，AI 就不会在无关紧要的内容上浪费篇幅，而是聚焦在你真正关心的问题上。

\section{Active-Prompt}

Active-Prompt 的核心想法是：先让 AI 回答一批问题，找出它回答得不一致或者不确定的那些（这些就是它的「薄弱环节」），然后针对这些难题，人工补充示例和注释，再让 AI 重新回答。

\textit{类比：先让实习生做一套模拟题，看他哪里老错，然后专门针对错题讲解。}

这个技巧的操作流程：

\begin{enumerate}
  \item 给 AI 一批问题，让它回答
  \item 检查哪些回答不够好或者前后不一致
  \item 针对这些问题，提供详细的示例和解释
  \item 把补充后的提示词再给 AI，让它重新回答
\end{enumerate}

来看一个实际场景。假设你需要 AI 帮你做客户反馈的情感分类。你先让它分类一批反馈，发现大部分「正面」和「负面」的分类都很准确，但碰到带有反讽语气的评论（比如「哦，你们的售后真是太棒了呢」），AI 总是判断错误。

这时候，你就可以针对这类「反讽」场景，补充几个带标注的示例：

\begin{promptbox}
以下是一些客户反馈的情感分类示例，请特别注意带有反讽或讽刺语气的表述：

反馈：「产品质量非常好，用了一年都没问题。」
分类：正面

反馈：「哦，你们的售后真是太棒了呢，打了三天电话都没人接。」
分类：负面（反讽）

反馈：「谢谢你们的快递，三周才到，我还以为包裹去环游世界了。」
分类：负面（反讽）

现在请对以下反馈进行情感分类：
（粘贴新的反馈内容）
\end{promptbox}

通过这种「先诊断弱点，再针对性补强」的方式，你可以用最小的工作量获得最大的质量提升。这个方法特别适合需要批量处理类似任务的场景，比如数据标注、分类文档，或者批量生成特定格式的内容。

\section{检索增强生成（RAG）}
\label{sec:rag}

检索增强生成（Retrieval-Augmented Generation，简称 RAG）是目前非常热门的一个技术。它的核心思路是：让 AI 在回答问题之前，先从你提供的文档或数据库里检索相关内容，然后基于检索到的内容来回答。

\textit{类比：不让实习生凭记忆回答客户问题，而是要求他先查公司知识库，找到依据再回答。}

为什么需要 RAG？因为 AI 的训练数据有截止日期，它不知道最新的信息；它也不了解你公司内部的规章制度、产品文档等私有知识。通过 RAG，你可以让 AI 「看到」这些内容，大幅减少它胡说八道的概率。

\textbf{简单的 RAG（手动版）：}

你可以直接在提示词里附上参考资料：

\begin{promptbox}
请根据以下产品文档，回答客户的问题。如果文档里没有提到的信息，请直接说「这个问题我需要确认后再回复您」。

【产品文档内容】

（粘贴你的文档内容）

客户问题：你们的退货政策是什么？
\end{promptbox}

\textbf{自动化的 RAG：}

更高级的大型 RAG 工程系统：利用 Embedding 模型把你的文档切成数百个段落碎块（Chunks）并转化为高维向量数组存入向量数据库（Vector DB），当用户提问时，系统会先启动距离相似度计算（如余弦相似度）自动检索最匹配的 TOP 5 文档块，再把它们打包当成背景 Context 发射给大模型融合宣读。市面上热门的各类「企业级专属知识库」及各种智能问答机器人的核心底层基石清一色都是 RAG。

对于绝大多数没有代码基础的零基础用户而言，这种「手动粘贴附带文档资料」的新手版 RAG，就已经能在工作流中产生降维打击般的威力了。你需要永远信奉一个原则：\textbf{与其让模型全凭它的赛博记忆去解答问题，不如直接把相关说明书「喂」进它的上下文里，这也是避免 AI 职场翻车的最强保险丝。}

这六种进阶技巧的共同目标都是一样的：通过更精准的引导方式，让 AI 给出确定性更强、更可靠的回答。不需要一次性全部记住，先留个印象，等在实践中遇到了瓶颈，再回到这一章查阅对应的技巧即可。如果你好奇提示工程还能如何进一步拓展，下一章将介绍一些更前沿的高阶方法。

%!%
% Copyright (c) 2026 mingcheng <mingcheng@outlook.com>
% 
% File: ch08-prompting-advanced.tex
% Author: mingcheng <mingcheng@outlook.com>
% File Created: 2026-03-27 15:50:02
% 
% Modified By: mingcheng <mingcheng@outlook.com>
% Last Modified: 2026-03-29 07:50:31
%%

% !TeX root = ../prompt-engineering-guide-for-beginners.tex

\chapter{提示工程高阶：前沿探索与自动化}
\label{ch:prompting-advanced}

这一章介绍的技巧相对前沿，有些涉及较复杂的概念。
如果你刚开始学习提示工程，完全可以先跳过这一章，等前面的基础和进阶技巧用熟了再回来看。
把这章当作一个「技术全景图」，了解一下提示工程还有哪些更高级的玩法。

\section{思维树（Tree of Thoughts, ToT）}

在第~\ref{ch:prompting-basics}~章我们学了链式思考（CoT），让 AI 沿着一条推理路径一步步思考。但有些问题不是一条直线就能解决的，中间可能有多个分叉点，每个分叉点都可以有不同的选择。

思维树就是让 AI 在每个关键节点展开多个分支，在这个过程中采用广度优先搜索（BFS）或深度优先搜索（DFS）的算法思路评估哪条路最有希望，然后继续深入探索。如果某条路走不通（死胡同），就执行回溯（Backtracking）到上一个节点，换一条路试试。

\textit{类比：下棋时在脑子里推演好几步，每一步都想几种走法，选最优的继续。AI 也可以这样。}

\textbf{适用场景：}需要「多步试错」的规划类任务、创意类任务，或者有多个可行方案需要比较优劣的决策问题。

\textbf{使用方法：}

\begin{promptbox}
我需要为一家新开的奶茶店制定开业推广方案。

请按照以下方式思考：
\begin{enumerate}
  \item 先列出 3 种不同的推广策略方向
  \item 对每种策略，分析其优缺点和预期效果
  \item 选择最有潜力的策略，深入展开具体执行方案
  \item 如果最终方案存在明显风险，回到第 1 步考虑其他方向
\end{enumerate}
\end{promptbox}

\section{自动推理并使用工具（ART）}

现在越来越多的 AI 系统不只是「说」，还能「做」。ART 让 AI 在推理过程中自动判断：「这一步我需要查一下资料」「这个数据我应该用计算器算一下」，然后自动调用相应的工具。

\textit{类比：聪明的助理在写报告时，遇到不确定的数据会自己去查资料、用计算器验算，而不是全靠脑补。}

\textbf{实际例子：}当你让 ChatGPT（联网版）回答「今天上海的天气怎么样」时，它会自动判断需要搜索实时信息，然后调用搜索工具获取结果，再把搜索结果整合到回答中。

这种能力在很多现代 AI 工具中已经内置了。比如 ChatGPT 的代码解释器可以自动运行 Python 代码来做精确计算，Kimi 和 ChatGPT 的联网版可以自动搜索最新信息。你在使用这些工具时可能已经体验过了，只是没有意识到背后的机制就是 ART。

作为用户，你需要注意两件事：

\begin{itemize}
  \item \textbf{主动启用工具功能}：有些工具需要你手动开启联网搜索或代码执行功能。遇到需要实时信息或精确计算的任务时，记得检查一下相关功能有没有打开。
  \item \textbf{明确指示 AI 使用工具}：如果你发现 AI 在「心算」一个复杂的数学问题，可以直接告诉它「请写代码来计算，不要口算」，它通常会乖乖执行。同样，如果你需要最新数据，可以提示它「请搜索最新的信息来回答」。
\end{itemize}

简单记住一个判断标准：凡是涉及精确数字、实时信息或大规模数据处理的任务，让 AI 调用工具来完成永远比让它「脑补」靠谱。

\section{自动提示工程师（Automatic Prompt Engineer, APE）}

写提示词本身也是一项可以自动化的任务。APE 的思路是：让 AI 自己生成大量不同版本的提示词，然后在测试集上评估它们的效果，自动选出表现最好的那一版。

\textit{类比：你不自己调配咖啡豆的比例，而是让机器人自动试了 100 种配方，选出最好喝的那杯。}

这个技巧目前主要用在需要大批量处理任务的场景中（比如做数据标注、内容审核等），个人用户可能用不太到。但知道这个概念很重要，因为它揭示了提示工程领域的一个趋势：\textbf{未来，很多提示词的优化工作会由 AI 自己完成}。

其实你在日常使用中已经可以做一个简化版的 APE 了。当你对一条提示词的效果不满意时，可以直接让 AI 帮你改进：「这条提示词的效果不够好，请帮我优化它，使输出更加简洁和结构化。」让 AI 自己来优化提示词，往往能给出你意想不到的改进。

\section{ReAct 框架（Reasoning + Acting）}

ReAct 是「推理」（Reasoning）和「行动」（Acting）的结合。它让 AI 交替进行思考和行动：先想一想下一步该做什么，然后执行（比如搜索、查数据），根据执行结果再想一想下一步，反复循环直到完成任务。

\textit{类比：侦探破案时不是一口气推理到底，而是想一想，去现场找证据，再想一想，再去查资料……}

ReAct 的思维过程通常是这样的：

\begin{enumerate}
  \item \textbf{思考}：根据当前信息，我应该先做什么？
  \item \textbf{行动}：执行一个具体的操作（搜索、计算、查询等）
  \item \textbf{观察}：看看行动的结果是什么
  \item \textbf{思考}：基于新的结果，下一步该做什么？
  \item 重复上述过程，直到得出最终答案
\end{enumerate}

很多基于行业标准框架（如 LangChain、AutoGPT）打造的 AI 智能体（Agent）系统就是完全基于 ReAct 框架构建底层逻辑调度的。对于普通用户来说，了解这个框架有助于你更深刻地理解为什么有些拥有“外部工具调用能力（Function Calling）”的 AI 系统能脱离纯文本限制，通过执行 API 请求来自动完成一系列复杂的多步骤任务。

\section{Reflexion（自我反思）}

Reflexion 让 AI 在完成任务之后「回头看看」自己的回答，检查是否有遗漏、错误或不合理的地方，然后自行修正。

\textit{类比：实习生写完报告后自己先通读一遍，标出不确定的地方，修改后再交给你。}

你可以在提示词中要求 AI 进行自我反思：

\begin{promptbox}
请帮我写一封商务邮件。写完之后，请你自己检查一遍，看看是否有以下问题：
\begin{itemize}
  \item 语气是否恰当
  \item 是否有语法或拼写错误
  \item 信息是否完整
  \item 有没有可能被误解的表述
\end{itemize}
如果发现问题，请直接修正并给出最终版本。
\end{promptbox}

这个技巧简单但有效。在很多场景下，加一句「请检查你的回答是否有误」就能提升输出质量。

\section{程序辅助语言模型（PAL）}

AI 在自然语言推理方面很强，但在精确计算方面容易出错。PAL 的思路是：让 AI 遇到需要计算的部分时，不要用自然语言「心算」，而是写一段代码来执行。

\textit{类比：让实习生别心算了，直接打开 Excel 用公式算，结果靠谱得多。}

在支持代码执行的 AI 工具中（比如 ChatGPT 的代码解释器），你可以这样提示：

\begin{promptbox}
我有一组销售数据（见下表）。请用 Python 代码计算每个月的同比增长率，并找出增长最快的三个月。不要心算，请写代码运行。
\end{promptbox}

PAL 的关键在于一句话：\textbf{告诉 AI 别心算，用代码}。只要你在提示词里明确要求「写代码来计算」「用程序来处理」，AI 就会切换到程序化的工作模式，输出的结果会精确得多。

除了数学计算，PAL 在这些场景中也很好用：日期推算（计算两个日期之间间隔多少个工作日）、字符串处理（批量替换文本中的格式）、数据统计（从一堆数据中找出中位数、标准差等）。这些任务让 AI 用自然语言推理容易出错，但写成代码就是小菜一碟。

\section{多模态思维链提示（Multimodal CoT）}

传统的链式思考只处理文字，但现在的 AI 模型越来越多地支持图片、图表等视觉信息。多模态 CoT 就是把视觉信息也纳入推理过程：AI 不只是读文字，还能「看图」一步步分析。

\textit{类比：实习生不只是读文字报告，还能看着图表一步步分析趋势并给出结论。}

\textbf{示例：}

\begin{promptbox}
请看这张销售趋势图。

1. 先描述图表中显示的整体趋势
2. 找出明显的拐点或异常值
3. 分析可能的原因
4. 给出你的预测和建议
\end{promptbox}

\section{基于图的提示（Graph Prompting）}

有些问题涉及复杂的实体关系，比如人物关系、组织架构、知识网络等。基于图的提示让你用节点和连线的方式描述这些关系，帮助 AI 更好地理解和推理。

\textit{类比：画一张人物关系图，让实习生顺着关系线去回答「张三和李四是什么关系」。}

\textbf{示例：}

\begin{promptbox}
以下是一个公司的汇报关系：
\begin{itemize}
  \item 张总（CEO）管理王总（CTO）和李总（CFO）
  \item 王总（CTO）管理赵经理（研发）和孙经理（测试）
  \item 李总（CFO）管理周经理（财务）和吴经理（审计）
  \item 赵经理管理开发团队的陈工和刘工
\end{itemize}

请问：陈工的直属上级是谁？如果陈工想找 CEO 反映问题，中间需要经过几层汇报？
\end{promptbox}

\section{元提示（Meta-Prompting）}

到目前为止，我们写的提示词都是针对具体任务的。元提示的思路更高一层：用一个「总调度」提示来管理和协调多个子任务。

\textit{类比：你不直接干活，而是当总导演，安排策划组、设计组、文案组各自完成任务，最后你来审核拼装。}

\textbf{示例：}

\begin{promptbox}
你是一个任务调度器。我需要完成以下项目：为一款新上线的健身 App 制定上线推广方案。

请按照以下步骤进行：
\begin{enumerate}
  \item 先以「市场分析师」的身份，分析当前健身 App 市场的竞争格局
  \item 再以「品牌策划师」的身份，确定我们的差异化定位和核心卖点
  \item 然后以「新媒体运营」的身份，制定各渠道的具体推广计划
  \item 最后，汇总以上分析，输出一份完整的推广方案提纲
\end{enumerate}
\end{promptbox}

元提示非常适合处理需要多种专业视角协作的复杂任务。它本质上是把「角色扮演」和「提示链」结合起来，再加上一个统筹管理的层级。

\bigskip

以上就是提示工程的高阶技巧全景。这些技巧中有一些已经在主流 AI 工具中有了直接的应用，有一些还停留在学术研究阶段。但了解它们的存在，能帮你在遇到复杂问题时找到更多的思路和灵感。

\begin{tipbox}
这一章的绝大部分技巧属于「知道有这回事」就够了的层级。日常使用中，前两章的基础和进阶技巧已经能覆盖绝大多数场景。只有当你遇到普通方法解决不了的复杂问题时，再回来翻阅本章寻找灵感。
\end{tipbox}

方法论讲完了，接下来我们直接上干货。下一章会提供一批拿来即用的提示词模板，覆盖总结、翻译、写作、代码等高频场景。

%!%
% Copyright (c) 2026 mingcheng <mingcheng@outlook.com>
% 
% File: ch09-prompt-examples.tex
% Author: mingcheng <mingcheng@outlook.com>
% File Created: 2026-03-27 15:50:02
% 
% Modified By: mingcheng <mingcheng@outlook.com>
% Last Modified: 2026-03-29 07:50:36
%%

% !TeX root = ../prompt-engineering-guide-for-beginners.tex

\chapter{提示词示例：拿来即用的场景模板}
\label{ch:prompt-examples}

前面几章讲了很多方法论，这一章直接上「干货」。
针对日常工作和生活中最常见的场景，我准备了一批可以直接复制使用的提示词模板。
你只需要把里面的具体内容替换成你自己的需求，就能拿来即用。

\section{文本概括与总结}

日常工作中最常见的场景之一就是「读了一大堆东西，但时间不够用」。无论是会议纪要、行业报告还是长篇文章，AI 都能帮你快速抓住重点。以下几个模板覆盖了最常见的总结需求。

\subsection*{长文章一句话摘要}

\begin{promptbox}
请用一句话概括以下文章的核心观点，要求简洁精炼，不超过 50 个字：

\verb|"""|

（在这里粘贴你的文章内容）

\verb|"""|
\end{promptbox}

\subsection*{会议记录提炼行动项}

\begin{promptbox}
以下是今天的会议记录。请帮我提炼出所有的待办事项（Action Items），以编号列表形式列出，每条包含：负责人、具体任务、截止时间。如果会议记录中没有明确负责人或截止时间，请标注为「待确认」。

\verb|"""|

（粘贴会议记录）

\verb|"""|
\end{promptbox}

\subsection*{论文/报告核心观点提取}

\begin{promptbox}
请阅读以下报告，提取 3 到 5 个核心观点。每个观点用一句话总结，并标注出自报告的哪个部分。最后用 2 到 3 句话写一段整体摘要。

\verb|"""|

（粘贴报告内容）

\verb|"""|
\end{promptbox}

\section{信息提取与分类}

面对一堆杂乱的原始数据或文本，手动整理既枯燥又容易出错。AI 在这类「从非结构化文本中提取结构化信息」的任务上表现非常出色，几秒钟就能完成你手工要做半小时的分类和提取工作。

\subsection*{从产品评论中提取情感倾向}

\begin{promptbox}
请分析以下产品评论，将每条评论的情感分类为「正面」「负面」或「中性」，并用一句话说明判断理由。以表格形式输出，列名为：评论内容、情感分类、判断理由。

\verb|"""|

1. 这个产品质量太棒了，用了半年一点问题都没有！
2. 包装挺好看的，但实际用起来和描述不太一样。
3. 还行吧，没什么特别的。
4. 第二天就坏了，客服还不给退，差评！

\verb|"""|
\end{promptbox}

\subsection*{简历关键信息结构化提取}

\begin{promptbox}
请从以下简历中提取关键信息，以 JSON 格式输出，字段包括：姓名、联系方式、最高学历、毕业院校、工作年限、最近一份工作的职位和公司名称、核心技能（列表）。

\verb|"""|

（粘贴简历内容）

\verb|"""|
\end{promptbox}

\subsection*{客服工单自动分类}

\begin{promptbox}
以下是一批客服工单的问题摘要。请将它们自动分类到以下类别中：「产品质量」「物流配送」「退换货」「账户问题」「咨询建议」「其他」。以表格形式输出，列名为：工单编号、问题摘要、分类。

\verb|"""|

1. 收到的商品有破损
2. 订单显示已签收但我没收到
3. 想问一下你们的会员积分怎么用
4. 密码忘了，找回密码的邮件收不到
5. 买了两件，想退一件

\verb|"""|
\end{promptbox}

\section{文本转换与翻译}

翻译和改写是 AI 的传统强项。但「翻译」远不止是中英互译，它还包括在不同语言风格之间的转换：把学术语言变成大白话，把代码变成注释，把口语变成书面语。下面几个模板帮你搞定这些场景。

\subsection*{中英互译并保持语气}

\begin{promptbox}
请将以下中文翻译成英文。要求：保持原文的语气和风格（正式/口语化/幽默等），不要逐字直译，确保英文表达自然流畅。

\verb|"""|

（粘贴中文内容）

\verb|"""|
\end{promptbox}

\subsection*{把学术论文改写成科普文章}

\begin{promptbox}
以下是一篇学术论文的摘要。请将它改写成一篇面向普通读者的科普短文，要求：不使用专业术语（如果必须使用则加括号解释），语气轻松易懂，控制在 300 字以内。

\verb|"""|

（粘贴论文摘要）

\verb|"""|
\end{promptbox}

\subsection*{代码注释自动生成（代码辅助解析）}

\begin{promptbox}
请为以下代码添加详细的中文注释。要求：每个函数加上功能说明的文档注释，关键逻辑处加上行内注释，注释用中文书写，语言简洁明了。重点指出潜在的性能瓶颈或安全漏洞。

\verb|"""|

（粘贴代码）

\verb|"""|
\end{promptbox}

\begin{tipbox}
在进行代码或文本结构的重构、翻译时，如果你只需要复制最终的输出结果而不需要看 AI 啰里啰嗦的开场白（例如“好的，这是您的代码：”），建议在 Prompt 末尾加上一行「请直接输出结果本身，不需要任何多余的解释与寒暄」。这将极大提高复制粘贴的效率。
\end{tipbox}

\section{职场生存篇}

写周报、发邮件、做汇报……这些职场日常看似简单，却经常消耗你大量的时间和精力。用好以下模板，你可以把这些「文字活」的时间压缩到原来的几分之一，把精力留给更重要的事情。

\subsection*{一键生成周报}

\begin{promptbox}
你是一名经验丰富的职场人士。请根据以下工作要点，帮我写一份本周工作周报。格式要求：分为「本周完成」「进行中」「下周计划」三个板块，每个板块用编号列表，语气专业简洁。

本周做了这些事：
\begin{itemize}
  \item （列出你本周做的事情）
\end{itemize}
\end{promptbox}

\subsection*{润色年终总结}

\begin{promptbox}
请帮我润色以下年终工作总结。要求：保持原意不变，提升文字的专业度和可读性，语气沉稳大气，适当加入量化的表述来增强说服力。如果原文中有口语化的表达，请替换为更书面的说法。

\verb|"""|

（粘贴你的年终总结草稿）

\verb|"""|
\end{promptbox}

\subsection*{写一封得体的催款/拒绝邮件}

\textbf{催款邮件：}
\begin{promptbox}
请帮我写一封催款邮件。背景：对方是合作了两年的老客户，有一笔 5 万元的款项已经逾期 45 天。要求：语气礼貌但坚定，既要维护客户关系又要明确表达催款诉求，提及合同约定的付款周期，控制在 200 字以内。
\end{promptbox}

\textbf{婉拒邮件：}
\begin{promptbox}
请帮我写一封婉拒邮件。背景：一位同行邀请我参加下周的行业论坛做分享嘉宾，但我确实没有时间。要求：语气真诚表达感谢和歉意，给出合理的理由，留有未来合作的余地，控制在 150 字以内。
\end{promptbox}

\subsection*{长文档秒变大纲}

\begin{promptbox}
请阅读以下文档，将其内容整理成一份多层级的大纲。要求：一级标题不超过 7 个，每个一级标题下用要点列出核心内容，整体结构清晰、层次分明。

\verb|"""|

（粘贴文档内容）

\verb|"""|
\end{promptbox}

\section{内容创作篇}

如果你需要做自媒体运营、写营销文案或者制作短视频，AI 可以成为你最高效的「文案搭档」。它不一定能直接产出完美的成品，但绝对能帮你快速生成一个高质量的起点，让你在此基础上修改润色。

\subsection*{小红书种草文案}

\begin{promptbox}
你是一位小红书爆款文案写手。请为以下产品写一篇种草笔记：

产品：（产品名称和简介）

要求：
\begin{itemize}
  \item 标题要有吸引力，包含 emoji
  \item 正文分段清晰，每段不超过 3 行
  \item 要有真实感和个人体验感，不要太像广告
  \item 适当使用 emoji 增加阅读趣味
  \item 结尾加上 3 到 5 个相关话题标签
  \item 总字数 200 到 300 字
\end{itemize}
\end{promptbox}

\subsection*{短视频脚本规划}

\begin{promptbox}
请帮我写一个 60 秒短视频脚本。

主题：（你的视频主题）

要求：
\begin{itemize}
  \item 分为：开头hook（5秒抓住注意力）、正文内容（45秒）、结尾引导互动（10秒）
  \item 每个部分标注预计时长、画面描述、台词/旁白内容
  \item 语气口语化，适合真人出镜
  \item 开头第一句话要能让用户停下来不划走
\end{itemize}
\end{promptbox}

\subsection*{取标题的 5 种姿势}

\begin{promptbox}
请为以下文章内容提供 5 个不同风格的标题，并标注每个标题的风格类型：

1. 悬念型（引发好奇心）
2. 数字型（列出具体数量）
3. 痛点型（直击读者问题）
4. 对比型（制造反差感）
5. 利益型（突出读者收益）

文章内容概要：（简要描述你的文章内容）
\end{promptbox}

\section{生活与学习篇}

AI 不只是工作工具，在日常生活和学习中也能帮上大忙。从规划旅行到练习口语，从理解复杂概念到制定学习计划，以下模板覆盖了几个最实用的场景。

\subsection*{定制旅游攻略}

\begin{promptbox}
请帮我做一份旅游攻略。

目的地：（地点）

出行时间：（日期和天数）

预算：（总预算）

人数：（几个人、有没有老人小孩）

偏好：（美食/自然风光/历史文化/购物/其他）

要求：以日程表形式规划，每天包含上午、下午、晚上的安排，注明景点、交通方式、预计花费和餐饮推荐。
\end{promptbox}

\subsection*{AI 英语口语陪练}

\begin{promptbox}
你是一位耐心的英语口语老师。接下来我们进行一对一的英语对话练习。

规则：
\begin{itemize}
  \item 我用英语和你对话，你也用英语回复
  \item 如果我的表达有语法错误或不自然的地方，请在回复后用中文指出并给出更好的表达方式
  \item 用日常对话的语气，不要太学术
  \item 难度适中，适合中级英语学习者
\end{itemize}

我们从这个话题开始：今天的天气。请你先开始。
\end{promptbox}

\subsection*{复杂概念通俗化}

\begin{promptbox}
请用最通俗易懂的语言解释「量子力学」这个概念。要求：像讲给一个 10 岁小朋友听那样解释，可以用生活中的例子做类比，避免任何专业术语，控制在 200 字以内。
\end{promptbox}

\section{代码与技术篇}

即使你不是专业的程序员，AI 也能帮你写简单的脚本、排查代码错误、生成测试用例。下面的模板针对几个最常见的编程辅助场景，你只需要把自己的代码或需求填进去就能使用。

\subsection*{用自然语言生成代码}

\begin{promptbox}
请用 Python 写一个脚本，实现以下功能：

1. 读取当前目录下所有的 .csv 文件
2. 将每个文件中的「日期」列转换为统一的 YYYY-MM-DD 格式
3. 合并所有文件到一个总表
4. 按日期排序后输出为一个新的 merged.csv 文件

要求：添加必要的注释，使用 pandas 库，做好异常处理。
\end{promptbox}

\subsection*{AI 帮你 Debug}

\begin{promptbox}
以下 Python 代码运行时报错了，错误信息是：（粘贴错误信息）

请帮我：
1. 解释这个错误是什么意思
2. 找出代码中导致错误的具体位置
3. 给出修复后的完整代码
4. 解释修复的原因

\verb|"""|

（粘贴你的代码）

\verb|"""|
\end{promptbox}

\subsection*{自动生成单元测试}

\begin{promptbox}
请为以下 Python 函数编写单元测试。要求：使用 pytest 框架，覆盖正常输入、边界值和异常输入三种情况，每个测试用例加上清晰的注释说明测试目的。

\verb|"""|

（粘贴你的函数代码）

\verb|"""|
\end{promptbox}

\bigskip

以上模板只是一个起点。用的时候记得根据自己的实际需求调整细节，特别是背景信息、约束条件和输出格式。用得越多，你就越能感受到提示词的微小差异带来的巨大效果差别。

接下来的几章，我们会聊一些更实用的话题：多模态交互、翻车救援、安全与伦理，以及如何把零散的技巧整合成系统化的工作流。

% !TeX root = ../prompt-engineering-guide-for-beginners.tex

%%%%%%%%%%%%%%%%%%%%%%%%%%%%%%%%%%
%% 第十章：多模态提示
%%%%%%%%%%%%%%%%%%%%%%%%%%%%%%%%%%
\chapter{多模态交互：不只是打字聊天}
\label{ch:multimodal}

到目前为止，我们聊的都是文字文本层面的对话。但现在的顶级大语言模型进化极快，它们早已经不仅仅局限于会「读」和「写」了：它们还能看懂图片、透视复杂文档、听懂语音，甚至挥毫泼墨帮你画画。
这一章将带你领略这些被统称为「多模态（Multimodal）」的能力，以及如何使用专属的提示词点亮它们。

\section{看图说话：图片理解与分析}

越来越多的 AI 模型支持直接上传图片进行分析。你可以让 AI 做这些事：

\begin{itemize}
  \item \textbf{描述图片内容}：上传一张照片，让 AI 告诉你图里有什么
  \item \textbf{提取图片中的文字}（OCR）：拍一张名片、截一张表格，让 AI 把里面的文字提取出来
  \item \textbf{分析图表数据}：上传柱状图、折线图、饼图等，让 AI 帮你解读数据趋势
  \item \textbf{辅助设计}：上传一张界面截图，让 AI 分析布局并提出改进建议
\end{itemize}

写图片分析提示词的技巧：

\begin{promptbox}
请看这张图片并完成以下任务：
\begin{enumerate}
  \item 描述图片中的主要内容
  \item 提取图片中所有可见的文字信息
  \item 如果是图表，请分析数据并给出 3 个关键发现
\end{enumerate}
\end{promptbox}

关键要点是：\textbf{千万不要偷懒只说一句「看看这张图」}。就像文字提示词一样，你需要明确界定你希望 AI 扮演什么角色、从图片中抓取什么信息、基于图片做什么具体分析。你的指令越有指向性，它的“视力”与“见解”就越敏锐。

\section{喂文档：让 AI 闪电阅读你的文件}

许多优秀的 AI 工具现在已原生支持直接上传外部文件（包括 PDF、Word、Excel、PPT 乃至长篇代码文件等），并可直接让 AI 阅读、交叉比对和总结回答其中的问题。这比你自己吭哧吭哧开文档去复制粘贴要优雅无数倍。

常见的文档处理神仙级应用场景：

\begin{itemize}
  \item \textbf{快速总结}：「请总结这份 PDF 的核心观点，不超过 300 字。」
  \item \textbf{文档问答}：「根据这份合同内容，乙方的违约责任具体是什么？」
  \item \textbf{多文档对比}：「请对比这两份报价单，找出价格差异超过 10\% 的项目。」
  \item \textbf{数据处理}：上传 Excel 文件，让 AI 帮你做统计、生成图表描述、找出异常值
\end{itemize}

使用文档功能时的提示词技巧：

\begin{promptbox}
我上传了一份公司年度财务报告（PDF）。请完成以下任务：
\begin{enumerate}
  \item 提取报告中的关键财务指标（营收、净利润、毛利率）
  \item 与上一年度进行对比（如果报告中有提及）
  \item 用 3 到 5 句话总结公司的财务健康状况
  \item 如果你发现报告中有值得注意的风险点，请单独列出
\end{enumerate}
\end{promptbox}

\begin{tipbox}
受限于算力和接口限制，许多工具对单次上传的文件体量（MB 数和总页数）设有硬性天花板。遇到动辄大几百页的超级年报，为了保证它检索的准确度，强烈建议拆分成多个章节文档按需分别投喂。
\end{tipbox}

\section{语音交互：用嘴代替键盘}

很多 AI 工具已经支持语音输入和语音输出。你可以直接用说话的方式和 AI 对话，不用打字。

语音交互最大的价值在于「解放双手」。想象一下这些场景：你正在开车上班的路上，突然想到一个文案的灵感，但没法打字；或者你在做饭时想让 AI 帮你整理一下明天的会议议程；又或者你想边散步边和 AI 练习英语口语。这些场景下，语音交互都比打字方便得多。

语音输入能为你带来「口述式写作」的奇特爽感。你在手写文案时可能会长时间卡壳，但嘴巴说起来往往滔滔不绝。你大可直接对着手机 App 版的模型把脑内繁杂琐碎的想法全都宣泄出来，最后收尾补上一句「请把刚才这些散乱的想法帮我整理成一篇结构严谨的述职大纲」。这对不擅长打底稿的人来说，无异于降维打击。

使用语音交互的心法小窍门：

\begin{itemize}
  \item 语意尽量连贯完整。不要说半句话就长久停顿，AI 极可能会将这半句话立刻当成完整指令去执行，导致结果错位。
  \item 在人名与数字等极精密的专有名词上放慢咬字。例如读错「三十五十万」和「三千五十万」将会产生巨大的业务灾难。
  \item 万一嘴瓢说错了，无需终止重来。你可以直接像和真人聊天一样补充纠正：「哎呀不对，我上面说的不是 X 而是 Y，请按 Y 作为标准」。AI 的语义容错理解能力足以完美剥离出正确的意图。
  \item 有节奏的留白。在表达完一个小节要点后维持一秒短暂的空隙，能大幅增强内置语音识别系统抓取断句标点符号的准确性。
\end{itemize}

\section{AI 生图：用文字画画}

现在有很多 AI 可以根据你的文字描述生成图片（叫做「文生图」），比如 DALL·E、Midjourney、Stable Diffusion 等。

写图片生成提示词和写文字提示词有一些不同，重点在于对视觉元素的描述。

一个好的生图提示词通常包含以下要素。你可以把它想象成给画家写的一张「创作需求单」，写得越具体，画家就越知道你想要什么：

\begin{itemize}
  \item \textbf{主体}：你想画什么？（一只猫、一座城市、一碗面）这是整个描述的基础，必不可少。
  \item \textbf{风格}：什么艺术风格？（水彩、油画、赛博朋克、吉卜力动画风格 、极简线条）这决定了图片的整体视觉感觉。
  \item \textbf{构图}：从什么角度看？（俯视、特写、全景、侧面）同一个对象从不同角度拍摄，给人的感觉完全不同。
  \item \textbf{光线和氛围}：什么样的光线和情绪？（暖色调、日落光线、雨天、温馨）光线和氛围往往是一张图从「还行」到「惊艳」的关键差别。
  \item \textbf{细节补充}：其他重要的视觉细节（背景模糊、高清、4K 等）
\end{itemize}

\textbf{示例：}

\begin{promptbox}
\textbf{简单版}：一只橘猫坐在窗台上

\textbf{详细版}：一只毛茸茸的橘猫坐在旧木窗台上，窗外是下雨的城市街景，画面风格为水彩插画，暖色调，柔和的自然光从窗户照进来，给人温馨治愈的感觉
\end{promptbox}

\begin{tipbox}
描述要素写得越详尽、越可视化，最终生成的图就越接近你的脑内设想。但请务必避免塞入太多「彼此互斥」的情境设定（例如要求「背景是正午的烈日」同时又要「充满昏暗阴郁的深夜惊悚感」），这会导致模型出现逻辑错乱甚至画出“克苏鲁”风格的缝合怪。
\end{tipbox}

\section{多模态提示词的特别红线}

在使用多模态功能时，有几点需要特别注意：

\begin{itemize}
  \item \textbf{模型能力差异}：不是所有 AI 都支持图片理解或文件上传。使用前确认你的工具支持哪些模态。
  \item \textbf{文件大小限制}：大多数工具对上传文件的大小有限制（通常几十 MB 以内），超大文件可能需要压缩或分割。
  \item \textbf{隐私保护}：上传图片和文档时要特别注意隐私。不要上传包含敏感个人信息、公司机密或他人隐私的内容。
  \item \textbf{图片中的文字}：AI 识别图片中文字的能力在不断提升，但仍然可能出错，特别是手写体、模糊图片或排版复杂的场景。重要信息建议人工核实。
  \item \textbf{生图的局限}：AI 生成的图片在细节上可能有瑕疵（比如手指数量不对、文字变形等），如果用于正式场合，需要仔细检查。
\end{itemize}

多模态能力是 AI 发展最快的方向之一。现在还做不到的事情，过几个月可能就能做了。保持对新功能的关注，适时尝试，你会发现越来越多意想不到的用法。

%!%
% Copyright (c) 2026 mingcheng <mingcheng@outlook.com>
% 
% File: ch11-troubleshooting.tex
% Author: mingcheng <mingcheng@outlook.com>
% File Created: 2026-03-27 14:57:14
% 
% Modified By: mingcheng <mingcheng@outlook.com>
% Last Modified: 2026-03-29 07:50:49
%%

% !TeX root = ../prompt-engineering-guide-for-beginners.tex

\chapter{翻车救援指南：模型不听话怎么办？}
\label{ch:troubleshooting}

在驱使 AI 当牛做马的过程中，你大概率会遇上各种各样的「翻车名场面」：回答得牛头不对马嘴、大张旗鼓地捏造数据、产出的文案像裹脚布一样又臭又长且毫无灵魂……
请安抚你的情绪，这些病症几乎都有特效解药。本章将手把手带你建立急救应对包，帮你逐一排除雷区。

\section{AI 开始一本正经地胡说八道（幻觉）怎么办？}

大模型最常见的致命劣根性被称为「幻觉」（Hallucination）。它的表现症状是：用极为扎实且自信的语气，堂而皇之地向你编造完全不存在的人名、事件、学术名言乃至莫须有的财务数据。

为什么会这样？因为 AI 的本质是「生成看起来合理的文字」，而不是「查找准确的事实」。当它不知道答案时，它不会说「我不知道」，而是会「编一个」。

\textbf{应对策略：}

\begin{itemize}
  \item \textbf{在提示词中设定边界}：「请仅根据我提供的文档内容来回答。如果文档中没有相关信息，请直接说『我无法确认』，不要猜测或补充。」
  \item \textbf{用文档替代记忆}：把你手头的权威资料、操作手册或内部文件直接粘贴到对话中，让 AI 基于这些有据可查的内容来回答。
  \item \textbf{要求标注信息来源}：「请在你的结论中标注每个观点的依据来源（来自哪个机构、哪一年的数据等）。」这虽然不能完全杜绝捏造，但能帮你快速核实真伪。
  \item \textbf{重要信息务必人工核实}：涉及公司宣传、法律条款、财务数据等零容错的场景，一定要用搜索引擎或其他可靠来源进行交叉验证，不要直接采信 AI 的输出。
\end{itemize}

\section{回答太长、太空洞、太像「机器人」怎么办？}

有时候 AI 为了显得全面，会输出一大段字数不少但缺乏实际价值的内容，语言里充斥着「首先、其次、总而言之、希望对你有所帮助」等套话，读起来像是在做填充题。

\textbf{解决办法：}

\begin{itemize}
  \item \textbf{设定字数限制}：「回答控制在 100 字以内」「用不超过三个要点来回答」。
  \item \textbf{指定你想要的文风}：「用知乎高赞回答的风格来写，简洁有力，带点实际经验」。
  \item \textbf{禁止套话}：「不要在开头和结尾加任何客套话或总结性废话，直接给出核心内容」。
  \item \textbf{要求可执行的干货}：「跳过背景介绍，直接给我三条明天就能落地执行的具体建议」。
\end{itemize}

\begin{tipbox}
常备的救场句式：「你这版回答太冗长了，帮我浓缩成核心观点，去掉开场白和结尾的客套话。」
\end{tipbox}

来看一个具体的对比。假设你想让 AI 给你写一封拒绝加班的回复，如果不加任何约束：

\begin{promptbox}
帮我写一条拒绝加班的消息发给领导。
\end{promptbox}

AI 可能会给你一大段客客气气、反复表决心的冗长文字。现在加上约束：

\begin{promptbox}
帮我写一条拒绝加班的消息发给领导。要求：不超过 50 字，语气礼貌但坚定，不要解释原因，直接表明今晚有事无法加班。
\end{promptbox}

\begin{aibox}
王总您好，今晚已有无法调整的安排，加班恐怕难以配合，抱歉给您添麻烦。如明日有需要我可以优先处理。
\end{aibox}

同样的需求，加上字数、语气和内容的约束后，输出就从「裹脚布」变成了可以直接发出去的消息。

\section{挤牙膏式追问法}

AI 第一次的回答不够好是很正常的。好消息是，你可以通过多轮对话不断追问和微调，逐步「逼」出你想要的结果。

追问的关键在于：告诉 AI 哪里不满意，以及你想要什么样的改变。不要只说「不好」「重新写」，而是说清楚「哪里不好」「怎样才算好」。想象一下你在指导一个实习生修改文档的场景，你不会只说「这个不行」，而会说「第二段的语气太生硬了，改成更口语化的表达」。对 AI 也是一样的道理。

\textbf{常用的追问句式：}

\begin{itemize}
  \item 「这个方向不错，但请在第二点上再展开一些，加入具体案例。」
  \item 「语气太正式了，请改成更轻松的风格。」
  \item 「你漏掉了XX方面，请补充进去。」
  \item 「这个回答太理论化了，请用一个实际工作中的例子来说明。」
  \item 「请把上面的内容重新组织一下，按照重要性从高到低排列。」
  \item 「不够具体。请给出可以直接执行的步骤，而不是泛泛的建议。」
\end{itemize}

这些追问并不需要从头重写提示词，只是在现有对话的基础上继续调整。通常经过二到三轮追问，你就能拿到相当满意的结果了。如果追问了五六轮还是不行，那往往说明你的初始提示词需要重新设计，而不是继续「挤牙膏」。

\section{Prompt 迭代优化流程}

写提示词不是一次性的事情，而是一个不断迭代优化的过程。就像写代码需要调试（Debug）一样，好的提示词也是改出来的：

\begin{enumerate}
  \item \textbf{起草（Draft）}：先把核心需求和原始材料直接写进提示词中，不必纠结措辞，先发出去试水。
  \item \textbf{观察和诊断（Inspect）}：看 AI 返回的结果，找出问题出在哪里。是理解有偏差？段落太散？还是风格不对？
  \item \textbf{针对性修改（Iterate）}：回到提示词中，利用前面章节学到的技巧，针对具体问题追加约束或补充说明，然后让它重新生成。
  \item \textbf{验证和打磨（Verify）}：如果还有瑕疵，继续「写 $\to$ 看效果 $\to$ 改」这个循环，直到满意为止。
\end{enumerate}

\begin{tipbox}
极其高亮的一条实验铁律：\textbf{你如果在进行调优，请保证绝不贪刀，尽量每次仅为提示词修正并验证「一个维度的变量因素」}（不要一拍脑门同时把语气、字数和格式全加上）。当你同时混杂着做四五个颠覆性的指令改动时，万一产出品质全面雪崩暴走，你作为控制者根本无从排查是被哪一个草率的参数调整拖垮了整盘棋局。
\end{tipbox}

另外一个好习惯是：\textbf{为自己的顺手武器库打版}。当你在经过七八轮焦灼的推拉后，终于在这个极其复杂的场景榨取到了想要的那条满状态金牌提示词大纲，别忘了把它放进你的便签本供起来（就像程序员积攒极佳的脚手架代码那样），下次直接调用。这也是后续章节「搭建私有指令库」的底层逻辑。

\section{AI 的回答自相矛盾？}

你可能会发现：同样的问题，第一次问和第二次问的答案不一样。这不是 Bug，而是 AI 的正常行为。

原因在于第~\ref{ch:basic-concepts}~章提到的 Temperature 参数（参见第~\ref{sec:temperature}~节）。AI 的回答本质上是基于概率的「采样」，每次采样的结果可能不同，就像掷骰子一样。

\textbf{应对方法：}

\begin{itemize}
  \item 如果你需要稳定一致的回答（比如数据分析、分类任务），在支持参数设置的工具中把 Temperature 调低
  \item 用「自我一致性」技巧（参见第~\ref{sec:self-consistency}~节），让 AI 回答多次，取最一致的结果
  \item 在提示词中明确要求：「请给出最确定的答案，不要给出模棱两可的回复」
  \item 对于重要决策，多问几次，看看答案是否一致。如果不一致，说明这个问题可能超出了 AI 能可靠回答的范围
\end{itemize}

\section{长文本被「截断」或「遗忘」了？}

当你输入很长的文本让 AI 处理，或者和 AI 聊了很多轮之后，你可能发现它开始「忘记」之前说过的内容。

这是因为上下文窗口有大小限制（参见第~\ref{sec:context-window}~节）。当总内容超过上限时，最早的内容就会被丢弃。

\textbf{应对方法：}

\begin{itemize}
  \item \textbf{分段处理长文本}：不要一次性把整本书贴进去。把长文本分成几个段落，每次处理一段，最后让 AI 汇总。
  \item \textbf{先总结再深入}：如果需要 AI 分析一篇很长的文章，可以先让它写一个总结，然后基于总结再深入讨论。
  \item \textbf{适时开新对话}：如果当前的话题已经结束，开一个新的对话窗口，把上下文空间腾出来。
  \item \textbf{重复关键信息}：如果你在对话进行中发现 AI 「忘了」某些重要设定，直接再说一遍就行。比如「请记住，你的角色是一位严格的编辑，请继续用这个身份回答我的问题。」
  \item \textbf{使用支持长上下文的模型}：如果你经常需要处理长文本，选择上下文窗口更大的模型（比如 Kimi 支持超长文本，Claude 也支持很长的上下文）。
\end{itemize}

总结一下这章的核心心法：AI 不是用一次就完美的工具，它更像是一个需要你不断调教和引导的合作伙伴。遇到问题不要放弃，试着调整你的提示词，大多数「翻车」都能救回来。

解决了「怎么用好」的问题，接下来我们还需要聊一个同样重要的话题：怎么「用得安全」。下一章会讨论使用 AI 时的隐私保护、事实核查和伦理边界。

% !TeX root = ../prompt-engineering-guide-for-beginners.tex

%%%%%%%%%%%%%%%%%%%%%%%%%%%%%%%%%%
%% 第十二章：安全与伦理
%%%%%%%%%%%%%%%%%%%%%%%%%%%%%%%%%%
\chapter{安全与伦理：负责任地使用 AI}
\label{ch:safety-ethics}

AI 是一个强大的工具，但强大也意味着需要谨慎使用。
这一章不是要吓唬你，而是帮你了解使用 AI 时的「红线」和「灰色地带」，让你用得安心、用得明白。

\section{隐私保护：哪些内容不该输入 AI}

很多人在使用 AI 时，会不假思索地把各种信息贴进去。比如把一封包含客户联系方式的邮件直接粘贴给 AI 帮忙润色，或者把公司内部的未公开财务报表发给 AI 做分析。这些操作看似方便，但你需要知道一个技术常识：\textbf{在默认设置下，你输入给公共 AI 模型（如免费版 ChatGPT、Claude）的所有 Prompt 和文件，都有可能被用作语料库，参与下一代模型的训练}。一旦被吸收进模型权重（Weights）中，这些机密信息就有可能在别人提问时被“吐”出来。

\textbf{绝对不能直接输入公共 AI 的内容清单（数据红线）：}

\begin{itemize}
  \item \textbf{个人敏感信息}：身份证号、银行卡号、密码、手机号、家庭住址等
  \item \textbf{公司机密}：内部财务数据、未发布的产品计划、客户名单、商业合同等
  \item \textbf{他人隐私}：同事的聊天记录、客户的个人信息、他人的医疗或法律记录
  \item \textbf{受保护的数据}：受保密协议（NDA）约束的内容、政府涉密信息
\end{itemize}

\textbf{企业级数据安全与脱敏建议：}

\begin{itemize}
  \item \textbf{建立本地脱敏机制（Data Masking）}：在将数据发送给云端 AI 之前，先将真实姓名、具体金额、核心项目代号替换为空占位或虚拟标签（如 \verb|[Client_A]|、\verb|[Amount_X]|），等云端推理并返回格式化文本后，再通过本地工具一键映射替换回来。这是最实用的隔离手段。
  \item \textbf{采用具备数据隔离（Zero Data Retention）承诺的服务}：使用正规云服务提供商的 API（如 OpenAI API、Azure OpenAI、AWS Bedrock）或购买认证的企业版（Enterprise/Team 版）。它们在用户协议中明确承诺：API 调用期间既不会缓存用户会话记录，也不会收集这批数据用于基础模型的持续训练。
  \item \textbf{部署本地私有化模型（Local LLMs）}：对于涉密或极高隐私敏感度的工作由于不能连通外网，最佳选择是依靠诸如 Llama 3、Qwen 这样高性能的开源模型，部署在本地电脑或内网机房进行纯离线推理。
  \item \textbf{一个简单的判断标准}：\textbf{绝不将任何你不敢直接发给陌生人或公开在社交平台上的机要信息贴进网页端免费版 AI 的对话框中}。
\end{itemize}

\begin{tipbox}
各大 AI 主流平台（如 ChatGPT 的 Data Controls 设置中）通常都配有“聊天记录与训练（Chat History \& Training）”或类似开关。使用前请务必确认已经关闭了“参与模型提升”服务，尽量从源头消除信息泄漏风险。
\end{tipbox}

\section{事实核查：AI 说的不一定是对的}

前面的章节多次提到了 AI 的「幻觉」问题。这里要再次强调：\textbf{AI 的回答不能作为唯一的事实依据}。这不是否认 AI 的价值，而是帮你在正确的场景中使用它。

以下场景中需要特别警惕：

\begin{itemize}
  \item \textbf{医疗健康}：AI 不是医生。它可以帮你了解症状的一般信息，但不能替代专业诊断。身体不舒服，请去医院。
  \item \textbf{法律咨询}：AI 不了解你所在地区的具体法律法规的最新修订。涉及法律决策时，请咨询专业律师。
  \item \textbf{财务投资}：AI 给的投资建议没有考虑到你的具体财务状况和风险承受能力，也不能预测市场走势。
  \item \textbf{学术引用}：AI 可能会编造论文标题、作者和期刊名。如果你需要引用文献，一定要手动验证每一条引用是否真实存在。
\end{itemize}

\textbf{核心原则（Human-in-the-loop）}：始终坚守“人类控制回路系统”的角色定位。把 AI 当作「不知疲倦的初级研究员」或「初稿的草拟者」，而不是「真理机器」或「全科专家」。它是一个概率推断模型而非精准数据库。只要是在可能带来经济损失、人身威胁、声誉风险的业务决策中，你永远应当作为「最后一环」承担核验与兜底。

\section{版权与原创性}

AI 生成的内容在版权方面目前还有很多不确定性。不同国家和地区的法律对此有不同的看法，而且法律还在不断演变中。这个领域目前还没有明确的全球统一标准，但有几个关键点值得你了解：

\begin{itemize}
  \item \textbf{AI 生成的内容不一定享有版权保护}：在一些法律体系中，只有人类创作的作品才能受到版权保护。纯粹由 AI 生成的内容可能不被视为「作品」。
  \item \textbf{AI 可能无意中复制训练数据}：虽然概率很低，但 AI 有可能生成与训练数据中某段内容非常相似的文字。在发布前最好检查一下。
  \item \textbf{标注 AI 辅助创作是一个好习惯}：如果你的内容有很大比例是由 AI 生成的，标注一下可以体现诚信，也能避免潜在的争议。
\end{itemize}

\textbf{创作权属归属建议}：与其将 AI 视作代笔工具，不如将其设定为“大脑外脑”或“头脑风暴伴侣”。保留对核心思想的原创性把控，让它草拟思路大纲、润色语法结构或寻找论据支撑，而文脉的节奏感和独特价值观点则交由人类自行创作。如果在商业化文本中使用了大范围 AI 原样生成的内容，适度注明“部分灵感由 AI 在监督下生成”既是职业操守的表现，也能从源头上斩断不必要的版权隐患。

\section{AI 偏见与公平性}

AI 模型是从大量的互联网数据中学习的，而互联网上的内容本身就包含各种偏见（性别偏见、种族偏见、文化偏见等）。这意味着 AI 的输出可能也会反映出这些偏见。

\textbf{需要特别注意的场景：}

\begin{itemize}
  \item \textbf{招聘和人事}：用 AI 筛选简历或撰写职位描述时，检查是否有无意中排斥某些群体的措辞。比如 AI 可能会在职位描述中写出「适合年轻人」这样的表述，无意中排除了年龄较大的候选人。
  \item \textbf{内容创作}：AI 生成的文案可能带有刻板印象。比如让它写一个「优秀工程师」的人物敬述，它可能默认写成男性；写「护士」的故事可能默认写成女性。这些隐含的偏见可能在不经意间传递给你的受众。
  \item \textbf{评估和决策}：不要让 AI 单独做出影响他人的重要决策。AI 可以提供分析和建议，但最终的判断应该由人来做
\end{itemize}

\textbf{系统性防偏见（De-biasing）策略}：所有基于大语料训练的算法必然带有互联网历史数据的思维定式（即所谓的统计学分布偏差），对于涉及人文社科和资源分配的任务，在使用 AI 输出之前，务必要从平权与多元维度人工多审视一遍。在系统提示词或要求中，可常备防止偏见的防御性约束：\verb|“在生成用词时，请确保多文化包容性与性别、年龄中立化，绝对不准带入涉及种族、职业、性别或地域色彩等任何刻板偏见，应当保证平权立场与客观分析视角。”|

\section{学术诚信与职场规范}

不同的学校和公司对 AI 使用有不同的规定。有的鼓励使用，有的限制使用，有的明确禁止在特定场景中使用。

\begin{itemize}
  \item \textbf{学生}：在写作业、论文或考试时，先了解你所在学校对 AI 使用的具体政策。很多学校要求你在使用 AI 辅助时必须声明，有些学校完全禁止在特定作业中使用 AI
  \item \textbf{职场人士}：了解公司对 AI 工具使用的规定。有些公司禁止员工向公共 AI 服务输入公司数据，有些公司提供了经过审批的内部 AI 工具
  \item \textbf{自由职业者}：如果你的客户要求原创内容，使用 AI 辅助创作时最好提前沟通，避免交付后产生纠纷
\end{itemize}

\textbf{总结}：AI 是工具，怎么用取决于你。了解规则、尊重边界、保持诚实，你就能在合规的前提下充分发挥 AI 的价值。毕竟，工具本身没有好坏，关键在于使用它的人。

%!%
% Copyright (c) 2026 mingcheng <mingcheng@outlook.com>
% 
% File: ch13-workflow.tex
% Author: mingcheng <mingcheng@outlook.com>
% File Created: 2026-03-27 15:50:02
% 
% Modified By: mingcheng <mingcheng@outlook.com>
% Last Modified: 2026-03-29 07:17:27
%%

% !TeX root = ../prompt-engineering-guide-for-beginners.tex

\chapter{从新手到高手：打造你的 AI 工作流}
\label{ch:workflow}

到这里，你已经掌握了提示工程的核心技巧。这一章要帮你从「会用 AI」升级到「善用 AI」，把零散的技巧整合成系统化的工作流。

\section{建立你的个人提示词库}

如果你用心写了一条效果很好的提示词，千万别用完就忘了。把它保存下来，下次遇到类似的任务直接复用。

\textbf{怎么建提示词库？}

不需要什么高级工具，一个文档或笔记本就够了。关键是做好分类和标注：

\begin{itemize}
  \item \textbf{分类标签}：按业务场景分类（如「邮件写作」「数据分析」「文案创作」），同时可以加上策略标签（如 \#零样本、\#链式思考），方便快速查找。
  \item \textbf{模板化设计}：不要把提示词写死，把需要替换的内容用方括号标出（如 \verb|[输入文本]|、\verb|[目标受众]|、\verb|[输出格式]|），下次使用时只需要填空就行。
  \item \textbf{记录效果与局限}：不仅记下好用的提示词，也记下它在什么情况下效果不好。这些经验比提示词本身更有价值。
  \item \textbf{版本记录}：像写代码一样记录每次优化的过程（如 v1 只有基本指令，v2 加入角色设定后效果更好），帮助你积累调优经验。
\end{itemize}

\begin{tipbox}
当提示词的积累达到一定量级时（超过 20 条高频词谱），请务必使用专门的知识管理工具如 Notion（含数据库视图）、Obsidian（含反向链接），甚至一些专门的 Github 提示词仓或第三方 Prompt SDK 服务，确保自己总能在一个统一搜索界面下，在 5 秒钟内复制并替换上相关的预设变量。
\end{tipbox}

\textbf{示例格式：}

\begin{promptbox}
\textbf{名称}：小红书种草文案

\textbf{场景}：为产品生成小红书风格的种草笔记

\textbf{提示词}：你是一位小红书爆款文案写手。请为以下产品写一篇种草笔记……（完整提示词）

\textbf{效果}：语气和格式都不错，偶尔需要手动调整 emoji 的使用密度

\textbf{版本}：v2（v1 没有限定字数，输出太长）
\end{promptbox}

积累到一定量之后，这个提示词库就是你的「AI 武器库」，能让你的工作效率实现质的飞跃。

\section{定制你的 AI 助手}

很多 AI 平台现在都提供了「自定义助手」的功能，让你可以创建一个预设好角色、知识和规则的专属 AI：

\begin{itemize}
  \item \textbf{ChatGPT 的 Custom GPTs}：你可以创建自己的 GPT，预设系统提示词、上传参考文档，甚至接入外部工具
  \item \textbf{Claude 的 Projects}：可以为特定项目创建独立的对话空间，上传相关文档作为知识库
  \item \textbf{Kimi 的智能体}：支持自定义角色和知识库，适合搭建中文场景下的专属助手
  \item \textbf{字节的 Coze}（扣子）：可以组合多种插件和工作流，搭建功能更复杂的 AI 应用
\end{itemize}

\textbf{定制 AI 助手的一些实用方向：}

\begin{itemize}
  \item \textbf{专属知识库问答}：上传你的行业研报、产品文档等资料，打造一位了解你业务的「专属分析师」。这种基于检索增强生成（RAG，参见第~\ref{sec:rag}~节）的方式能有效减少幻觉。
  \item \textbf{日常文书自动化}：打造一位帮你处理周报、会议纪要等固定格式文档的助手。通过预设好模板和示例，你只需输入关键信息就能得到规范的输出。
  \item \textbf{模拟演练}：创建一位扮演「刁钻客户」或「严格审稿人」的对抗角色，用来在正式提交前测试你的方案、文档是否经得起挑战。
  \item \textbf{工作流联动}：通过接入外部工具（如搜索、代码执行、消息推送等），让 AI 不仅能对话，还能实际执行操作，迈向智能体（Agent）架构。
\end{itemize}

\begin{tipbox}
定制 AI 助手的核心在于 System Prompt（系统提示词，参见第~\ref{sec:system-prompt}~节）。在这段设定指令中清晰定义助手的角色定位、回答流程和禁止行为，它的质量直接决定了助手的表现上限。
\end{tipbox}

\section{把 AI 融入你的工具链}

如果你目前使用 AI 的方式还停留在「复制 AI 的回答，粘贴到其他软件」，那还有很大的提效空间。更高效的做法是把 AI 当作工具链中的一个环节，和你日常使用的其他软件协作起来。

\textbf{常见的 AI + 工具组合：}

\begin{itemize}
  \item \textbf{AI + Excel/表格工具}：让 AI 帮你写复杂的 Excel 公式、生成数据分析的思路，或者直接处理你上传的表格数据。比如你不会写 VLOOKUP 公式，只需要告诉 AI「我有两张表，A 表有订单号和金额，B 表有订单号和客户名，请帮我写一个公式把客户名匹配到 A 表」，几秒钟就搞定了。
  \item \textbf{AI + 笔记工具}（Notion / 飞书文档 / Obsidian）：用 AI 快速生成初稿，再在笔记工具中整理和完善。很多笔记工具现在已经内置了 AI 功能，可以直接在编辑器里调用。
  \item \textbf{AI + 邮件}：用 AI 起草邮件，复制到邮箱客户端中微调后发送。特别适合需要用英文写邮件但又不太自信的场景。
  \item \textbf{AI + Python/代码}：让 AI 帮你写脚本来自动化重复性工作。即使你不会写代码也没关系，你只要用自然语言描述你要做的事情就行，比如「写一个脚本，把这个文件夹里所有 PDF 文件的文件名改成日期开头的格式」。
  \item \textbf{AI + 搜索引擎}：先用 AI 了解一个话题的基本框架和关键概念，再用搜索引擎查找具体数据和最新信息。AI 帮你建立认知框架，搜索引擎帮你验证和补充细节，两者配合效率最高。
  \item \textbf{AI + 设计工具}：先用 AI 生成文案和创意方向，再用 Canva、Figma 等设计工具落地视觉执行。你也可以让 AI 帮你写设计需求文档，然后把文档交给设计师或设计工具。
\end{itemize}

核心思路是：\textbf{让 AI 做它擅长的事情（生成初稿、处理文字、提供思路），让专业工具做它擅长的事情（精确计算、视觉设计、数据存储）}。不要试图让 AI 做所有事情，也不要忽略它能帮上忙的环节。两者结合，才能发挥最大价值。

\section{持续学习：跟上 AI 的进化速度}

AI 领域的变化速度非常快。新的模型、新的功能、新的技巧在不断涌现。今天的最佳实践，过几个月可能就有了更好的替代方案。

\textbf{如何保持学习：}

\begin{itemize}
  \item \textbf{定期尝试新工具}：每隔一段时间试试新出的 AI 工具或新功能。不需要深入研究，快速体验一下就能知道是否对你有用
  \item \textbf{关注几个可靠的信息源}：选择 2 到 3 个你信任的科技媒体或博主，定期浏览他们关于 AI 的更新
  \item \textbf{加入社区}：找一个 AI 爱好者的社区（微信群、论坛、Discord 等），和其他人交流使用心得
  \item \textbf{在工作流微调中实践（Do-by-Learning）}：不要为了学提示词技巧而硬刷。当你下一次在代码报错（Error Stack Trace）、大量网页去重整理（Data Processing）或是写重复信件时，暂停三秒反问：“能不能构建一段指令喂给模型批处理”，真实痛点中的磨砺胜过阅读万字干货。
  \item \textbf{无视参数量，聚焦逻辑拆解能力，关注大语言模型推理路线}：不必因没弄懂 O1 模型的底层原理（RLHF 机制）与最新参数量突破便感到自己落伍了！本质上，大模型在几年内依然是强大的“文字概率接龙机”，万金油法则便是锻炼自己将问题切分为极小的微元组件。
\end{itemize}

\begin{tipbox}
掌握结构化的提示词框架（就像写作中的「三段式」和编程中的「设计模式」）就足以应对 80\% 的工作挑战。至于模型底层的技术原理和 Token 优化等深入话题，留给专门研究这个方向的工程师就好。你的首要目标是让 AI 切实提升你的工作效率。
\end{tipbox}

记住，学习 AI 不是为了追赶潮流，而是为了让你的工作和生活更高效。找到适合自己的节奏，持续进步就好。

%!%
% Copyright (c) 2026 mingcheng <mingcheng@outlook.com>
% 
% File: conclusion.tex
% Author: mingcheng <mingcheng@outlook.com>
% File Created: 2026-03-28 09:28:17
% 
% Modified By: mingcheng <mingcheng@outlook.com>
% Last Modified: 2026-03-29 07:17:08
%%

% !TeX root = ../prompt-engineering-guide-for-beginners.tex

\chapter{结束语}

恭喜你读到了这里。

还记得第一章里那个比喻吗？去餐厅点菜，说「来个菜」和说「来一份不辣的番茄炒蛋，少油少盐」，结果天差地别。跟 AI 打交道的道理完全一样。一路读下来你会发现，所谓的提示工程，说到底就是和 AI 好好说话，把话说清楚这件事。角色设定、上下文补充、输出约束、链式推理，这些听起来各有各的名字，但本质上都是在帮你把一句模糊的需求拆成一段 AI 能准确执行的指令。

读完这本小册子只是起点。提示工程是一项需要在实践中打磨的手艺，就像学骑自行车一样，光看教程永远学不会，你得自己骑上去摔几次才能找到平衡感。明天写邮件的时候，整理会议纪要的时候，甚至帮孩子检查作业的时候，不妨打开 AI 试一试。第一次的结果大概率不完美，但没关系，换个角色、加点约束、给几个示例，再跑一遍，你就已经在做真正的提示工程了。每一次这样的迭代，都在把你对 AI 的理解往前推一步。

AI 领域的变化速度远超我们的想象，今天好用的技巧明天可能就有更好的替代方案。但你在这本小册子里学到的那些底层原则，「结构化表达、提供上下文、明确约束」，不会因为模型换代就失效。掌握了这些，新工具出来的时候你也能很快上手。而如果在使用过程中你踩到了有意思的坑，或者发现了一条特别好用的提示词，不妨分享给身边的朋友和同事。教是最好的学，在分享中你自己也会理解得更深。

最后，感谢你花时间读完这本小册子。希望它能成为你和 AI 打交道的一块敲门砖，帮你在工作和生活中真正用好这个时代最强大的工具。


\appendix

% !TeX root = ../prompt-engineering-guide-for-beginners.tex

%%%%%%%%%%%%%%%%%%%%%%%%%%%%%%%%%%
%% 第十四章：工具与资源
%%%%%%%%%%%%%%%%%%%%%%%%%%%%%%%%%%
\chapter{工具与资源}
\label{ch:resources}

这一章汇集了各种实用的工具、资源和速查表，方便你在需要的时候快速查阅。

\section{主流 AI 工具对比表}

\begin{center}
\begin{tabular}{p{2.2cm}p{2.5cm}p{2.8cm}p{2cm}p{3cm}}
  \toprule
  \textbf{工具} & \textbf{开发商} & \textbf{特点} & \textbf{免费额度} & \textbf{适用场景} \\
  \midrule
  ChatGPT & OpenAI & 功能全面，生态丰富 & 有（GPT-4o mini） & 通用，各类任务 \\
  Claude & Anthropic & 长文本处理能力强 & 有 & 写作、分析、编程 \\
  Gemini & Google & 深度整合谷歌生态 & 有 & 搜索、办公、多模态 \\
  Kimi & 月之暗面 & 超长文本，中文友好 & 有 & 中文场景，文档分析 \\
  DeepSeek & DeepSeek & 开源，推理能力强 & 有 & 推理、代码、通用 \\
  智谱 GLM & 智谱 AI & 清华系，API 丰富 & 有 & 开发集成 \\
  文心一言 & 百度 & 中文理解能力好 & 有 & 中文内容创作 \\
  通义千问 & 阿里 & 支持多模态 & 有 & 企业应用，代码 \\
  \bottomrule
\end{tabular}
\end{center}

\textit{说明：以上信息可能会随产品迭代而变化，请以各平台官网的最新信息为准。}

\section{提示词资源网站}

\begin{itemize}
  \item \textbf{Prompt Engineering Guide}（\url{https://www.promptingguide.ai/zh}）：系统全面的提示工程教程，支持中文，内容持续更新
  \item \textbf{FlowGPT}（\url{https://flowgpt.com}）：社区驱动的提示词分享平台，可以直接使用或收藏他人分享的提示词
  \item \textbf{PromptBase}（\url{https://promptbase.com}）：提示词交易平台，可以买卖高质量的提示词模板
  \item \textbf{Awesome ChatGPT Prompts}（GitHub）：开源的提示词集合，包含大量经过测试的角色扮演和任务型提示词
  \item \textbf{Learn Prompting}（\url{https://learnprompting.org}）：免费的提示工程在线课程，从入门到高级，有中文版
\end{itemize}

\chapter{常用高频提示词词汇表}
\label{ch:common-prompts}

在写提示词时，以下词汇可以帮助你更精确地表达需求：

\section{任务类动词}
\begin{center}
\begin{tabular}{llll}
  \toprule
  \textbf{中文} & \textbf{英文} & \textbf{中文} & \textbf{英文} \\
  \midrule
  总结 & Summarize & 分析 & Analyze \\
  解释 & Explain & 改写 & Rewrite \\
  翻译 & Translate & 比较 & Compare \\
  分类 & Classify & 提取 & Extract \\
  列举 & List & 评估 & Evaluate \\
  扩展 & Expand & 简化 & Simplify \\
  纠正 & Correct & 优化 & Optimize \\
  模拟 & Simulate & 预测 & Predict \\
  \bottomrule
\end{tabular}
\end{center}

\section{格式与风格描述词}
\begin{center}
\begin{tabular}{llll}
  \toprule
  \textbf{中文} & \textbf{英文} & \textbf{中文} & \textbf{英文} \\
  \midrule
  正式的 & Formal & 口语化的 & Conversational \\
  专业的 & Professional & 通俗易懂的 & Simple / Plain \\
  幽默的 & Humorous & 严肃的 & Serious \\
  简洁的 & Concise & 详细的 & Detailed \\
  客观的 & Objective & 有说服力的 & Persuasive \\
  \bottomrule
\end{tabular}
\end{center}

\section{输出格式关键词}
\begin{center}
\begin{tabular}{llll}
  \toprule
  \textbf{中文} & \textbf{英文} & \textbf{中文} & \textbf{英文} \\
  \midrule
  表格 & Table & 编号列表 & Numbered List \\
  要点列表 & Bullet Points & 思维导图 & Mind Map \\
  JSON 格式 & JSON Format & Markdown & Markdown \\
  分步骤 & Step by Step & 对比表 & Comparison Table \\
  \bottomrule
\end{tabular}
\end{center}

\chapter{提示工程术语表（中英对照）}
\label{ch:glossary}

\begin{center}
\begin{tabular}{p{4.5cm}p{8cm}}
  \toprule
  \textbf{术语} & \textbf{说明} \\
  \midrule
  提示词 / Prompt & 你输入给 AI 的指令或问题 \\
  提示工程 / Prompt Engineering & 设计和优化提示词的方法论 \\
  大语言模型 / LLM & Large Language Model 的缩写，AI 的核心技术 \\
  Token / 令牌 & 模型处理文本的基本单位 \\
  上下文窗口 / Context Window & 模型一次能处理的文本总量上限 \\
  Temperature / 温度 & 控制输出随机性的参数 \\
  幻觉 / Hallucination & AI 编造虚假信息的现象 \\
  零样本 / Zero-Shot & 不给示例，直接下指令 \\
  少样本 / Few-Shot & 通过提供几个示例来引导 AI \\
  链式思考 / Chain-of-Thought & 让 AI 展示推理过程 \\
  系统提示词 / System Prompt & 为整个对话设定规则的隐藏指令 \\
  RAG / 检索增强生成 & 先检索相关内容，再基于内容生成回答 \\
  微调 / Fine-Tuning & 用特定数据进一步训练模型 \\
  智能体 / Agent & 能自主执行多步骤任务的 AI 系统 \\
  多模态 / Multimodal & 同时处理文字、图片、语音等多种信息 \\
  推理模型 / Reasoning Model & 回答前先进行内部深度推理的模型（如 o1、DeepSeek R1） \\
  \bottomrule
\end{tabular}
\end{center}

\chapter{延伸阅读}
\label{ch:extended-reading}

如果你想进一步深入学习，以下是一些推荐的资源：

\section{在线课程}

\begin{itemize}
  \item \textbf{ChatGPT Prompt Engineering for Developers}（DeepLearning.AI）：由 Andrew Ng 和 OpenAI 联合推出的免费短课程，非常适合入门
  \item \textbf{Prompt Engineering for ChatGPT}（Coursera）：Vanderbilt University 提供的系统课程
\end{itemize}

\section{参考论文}

\begin{itemize}
  \item \textit{Chain-of-Thought Prompting Elicits Reasoning in Large Language Models}（Wei et al., 2022）
  \item \textit{Tree of Thoughts: Deliberate Problem Solving with Large Language Models}（Yao et al., 2023）
  \item \textit{ReAct: Synergizing Reasoning and Acting in Language Models}（Yao et al., 2022）
\end{itemize}

\section{社区}

\begin{itemize}
  \item Reddit: r/PromptEngineering、r/ChatGPT
  \item GitHub 上各种 Awesome 系列的提示词仓库
  \item 各大 AI 工具的官方社区和论坛
\end{itemize}

\chapter{常见问题 FAQ}
\label{ch:faq}

\section*{Q1：我应该选哪个 AI 工具？}
没有绝对的好坏，取决于你的需求。建议先免费试用几个，看哪个最顺手。中文用户可以优先试试 Kimi 或 DeepSeek，英文场景可以先试 ChatGPT 或 Claude。

\section*{Q2：免费版和付费版差别大吗？}
免费版足够日常学习和轻度使用。付费版通常提供更强的模型、更大的上下文窗口、更快的响应速度和更多的高级功能。如果你每天高频使用，付费版的体验提升很明显。

\section*{Q3：英文提示词是不是比中文效果好？}
在早期确实如此，因为模型的训练数据中英文占比更大。但现在主流模型的中文能力已经很强了，日常使用中文和英文的差距越来越小。用你更擅长的语言写提示词就好。

\section*{Q4：提示词有最佳长度吗？}
没有固定的最佳长度。简单的任务用一两句话就够了，复杂的任务可能需要几百字的详细说明。原则是：\textbf{说清楚就好，不要为了写长而写长，也不要为了简短而省略关键信息}。

\section*{Q5：AI 会取代我的工作吗？}
AI 更可能改变你的工作方式，而不是取代你。就像电脑和互联网的出现一样，它会淘汰一些岗位，但也会创造新的岗位。学会使用 AI 的人，会比不会使用的人更有竞争力。

\section*{Q6：我不懂技术，能学会提示工程吗？}
当然可以。提示工程的核心就是「把你的需求说清楚」，这是一项沟通能力，不需要任何硬核的代码技术背景。你已经读到了这本小册子的最后一章，说明你已经在路上了。

即便未来底层模型智能进一步提升，所谓的“自然语言编程”依旧是人类利用系统级思维与机器对话的桥梁。不管前端界面怎么变，逻辑拆解与约束能力永远是无法被替代的核心资产。

\section*{Q7：AI 的回答我能直接用吗？}
建议把 AI 的回答当作「初稿」而非「终稿」。在使用之前，花几分钟检查事实准确性、调整语气和细节，确保内容真正符合你的需求。

\section*{Q8：提示词有「标准答案」吗？}
没有。同一个需求可以有很多种写法，不同的模型对同一个提示词的反应也可能不同。找到适合你的场景和模型的写法，就是好的提示词。

\section*{Q9：把 AI 用好需要多长时间？}
基础技巧一两天就能上手。但真正熟练地使用 AI 来提高工作效率，需要在实际工作中不断练习和积累。保持好奇心，多尝试，你会越用越顺手。

\section*{Q10：AI 技术变化这么快，我学的东西会不会很快过时？}
模型和工具会不断更新，但提示工程的核心原则（说清楚需求、提供上下文、给出约束条件等）是不变的。掌握了这些原则，不管工具怎么变，你都能快速适应。

\begin{tipbox}
如果你有任何其他问题，或者想分享你的使用经验，可以在 GitHub 仓库提交 Issue。我会定期更新 FAQ 来解答大家的疑问。
\end{tipbox}

\chapter{更新记录}
\label{ch:changelog}

本小册子会随着 AI 技术的发展持续更新。以下是主要的版本变更记录：

\begin{center}
\begin{tabular}{p{2.5cm}p{2cm}p{8cm}}
  \toprule
  \textbf{日期} & \textbf{版本} & \textbf{更新内容} \\
  \midrule
  2026 年 3 月 & v1.1 & 优化内容和排版 \\
  2026 年 2 月 & v1.0 & 初版发布，包含完整内容 \\
  \bottomrule
\end{tabular}
\end{center}

这本小册子是开源的，欢迎大家参与更新和完善。如果你有任何建议或想贡献内容，可以在 GitHub 仓库提交 Pull Request。

\end{document}
